\documentclass[11pt,oneside]{article}

\usepackage[T1]{fontenc}
\usepackage{lmodern}
\usepackage{microtype}
\usepackage[margin=0.92in,headheight=15pt]{geometry}
\usepackage{amsmath,amssymb,amsthm,mathtools}
\usepackage{booktabs,longtable,tabularx,array,multirow}
\usepackage{enumitem}
\usepackage[dvipsnames]{xcolor}
\usepackage{xurl}
\usepackage{hyperref}
\usepackage{fancyhdr}
\usepackage{titlesec}
\usepackage{needspace}
\usepackage{listings}
\usepackage{seqsplit}

\definecolor{LeanBlue}{RGB}{28,86,140}
\definecolor{PaperGreen}{RGB}{20,112,72}
\definecolor{NewViolet}{RGB}{99,56,150}
\definecolor{ComputeOrange}{RGB}{166,86,20}
\definecolor{CautionRed}{RGB}{150,42,42}
\definecolor{MutedGray}{RGB}{88,96,105}
\definecolor{CodeBg}{RGB}{247,248,250}
\definecolor{RuleGray}{RGB}{95,103,112}

\hypersetup{
  colorlinks=true,
  linkcolor=LeanBlue,
  citecolor=PaperGreen,
  urlcolor=LeanBlue,
  pdftitle={Continuing the Alaoglu--Erdos Program},
  pdfauthor={OpenAI GPT-5.6 Pro; prepared for Vladimir Reshetnikov},
  pdfsubject={Synchronized Sturmian dynamics, effective logarithmic separation, and sharp p-adic determinant ceilings}
}
\urlstyle{same}

\pagestyle{fancy}
\fancyhf{}
\lhead{Alaoglu--Erd\H{o}s integer-exponent conjecture}
\rhead{Continuation report}
\cfoot{\thepage}
\renewcommand{\headrulewidth}{0.4pt}

\titleformat{\section}{\Large\bfseries}{\thesection}{0.72em}{}
\titleformat{\subsection}{\large\bfseries}{\thesubsection}{0.72em}{}
\titleformat{\subsubsection}{\normalsize\bfseries}{\thesubsubsection}{0.72em}{}

\setlist[itemize]{leftmargin=2em,itemsep=0.27em,topsep=0.38em}
\setlist[enumerate]{leftmargin=2.25em,itemsep=0.27em,topsep=0.38em}
\setlist[description]{leftmargin=3.7cm,labelsep=0.6cm,itemsep=0.35em,topsep=0.4em}
\setlength{\parindent}{0pt}
\setlength{\parskip}{0.53em}
\setlength{\emergencystretch}{3.5em}
\allowdisplaybreaks

\newtheorem{theorem}{Theorem}[section]
\newtheorem{proposition}[theorem]{Proposition}
\newtheorem{lemma}[theorem]{Lemma}
\newtheorem{corollary}[theorem]{Corollary}
\newtheorem{conjecture}[theorem]{Conjecture}
\newtheorem{question}[theorem]{Question}
\theoremstyle{definition}
\newtheorem{definition}[theorem]{Definition}
\newtheorem{observation}[theorem]{Observation}
\newtheorem{target}[theorem]{Research target}
\theoremstyle{remark}
\newtheorem{remark}[theorem]{Remark}
\newtheorem{warning}[theorem]{Caution}

\newcommand{\N}{\mathbb N}
\newcommand{\Nzero}{\mathbb N_0}
\newcommand{\Z}{\mathbb Z}
\newcommand{\Q}{\mathbb Q}
\newcommand{\R}{\mathbb R}
\newcommand{\C}{\mathbb C}
\newcommand{\Qbar}{\overline{\mathbb Q}}
\newcommand{\Sol}{\mathcal S}
\newcommand{\Cand}{\mathcal C}
\newcommand{\DeltaTri}{\mathcal T}
\newcommand{\fract}[1]{\left\{#1\right\}}
\newcommand{\distZ}[1]{\left\lVert #1\right\rVert_{\Z}}
\newcommand{\vpp}[2]{v_{#1}\!\left(#2\right)}
\newcommand{\supp}{\operatorname{supp}}
\newcommand{\Span}{\operatorname{span}}
\newcommand{\rank}{\operatorname{rank}}
\newcommand{\sha}[1]{\texttt{\seqsplit{#1}}}
\newcommand{\lean}[1]{\texttt{#1}}
\newcommand{\module}[1]{{\ttfamily\scriptsize\seqsplit{#1}}}
\newcommand{\statuslean}{\textcolor{LeanBlue}{\textbf{[Lean]}}}
\newcommand{\statusrepo}{\textcolor{LeanBlue}{\textbf{[repository build]}}}
\newcommand{\statusnew}{\textcolor{NewViolet}{\textbf{[new paper proof]}}}
\newcommand{\statusclassical}{\textcolor{MutedGray}{\textbf{[classical]}}}
\newcommand{\statuscomp}{\textcolor{ComputeOrange}{\textbf{[computation]}}}
\newcommand{\statusopen}{\textcolor{CautionRed}{\textbf{[open]}}}

\newenvironment{statusbox}[1]
{\begin{center}\begin{minipage}{0.95\textwidth}\raggedright\hrule\vspace{0.45em}\textbf{Status.} #1\par\vspace{0.35em}}
{\vspace{0.35em}\hrule\end{minipage}\end{center}}

\newenvironment{resultbox}
{\begin{center}\begin{minipage}{0.95\textwidth}\hrule\vspace{0.45em}}
{\vspace{0.35em}\hrule\end{minipage}\end{center}}

\lstdefinestyle{pythonstyle}{
  language=Python,
  basicstyle=\ttfamily\footnotesize,
  backgroundcolor=\color{CodeBg},
  frame=single,
  rulecolor=\color{RuleGray},
  numbers=left,
  numberstyle=\tiny\color{RuleGray},
  numbersep=8pt,
  showstringspaces=false,
  breaklines=true,
  columns=fullflexible,
  keepspaces=true,
  tabsize=4
}
\lstdefinestyle{leanstyle}{
  basicstyle=\ttfamily\footnotesize,
  backgroundcolor=\color{CodeBg},
  frame=single,
  rulecolor=\color{RuleGray},
  showstringspaces=false,
  breaklines=true,
  columns=fullflexible,
  keepspaces=true,
  tabsize=2
}

\begin{document}
\pagenumbering{gobble}

\begin{titlepage}
\centering
\vspace*{0.8cm}
{\Huge\bfseries Continuing the Alaoglu--Erd\H{o}s Program\par}
\vspace{0.55cm}
{\LARGE Synchronized Sturmian Dynamics, Effective Logarithmic Separation,\par}
\vspace{0.12cm}
{\LARGE and Sharp $p$-Adic Determinant Ceilings\par}
\vspace{0.9cm}
{\large A continuation of the attached unified report on\par}
\vspace{0.1cm}
{\large \url{https://github.com/VladimirReshetnikov/ProveIt}\par}
\vfill
\begin{minipage}{0.91\textwidth}
\small
\textbf{Verdict.}
I did not obtain an unconditional proof of the Alaoglu--Erd\H{o}s conjecture.  The exact
$2\times2$ logarithmic determinant remains the terminal obstruction.  The continuation does,
however, produce several rigorous advances and two fairly sharp no-go theorems for the most
natural determinant strategies.

A hypothetical primitive counterexample has a new exact description as one aperiodic
Sturmian word simultaneously driving a dyadic and a triadic rational recurrence.  Baker's
theory upgrades the elementary exponential continued-fraction bound in the source report to
an effective polynomial irrationality measure for the primitive generator.  On the
$p$-adic side, ordinary block Vandermonde valuations can be computed through all primes in the
fixed supports, giving a depth-sensitive asymptotic and a universal $87.1049\ldots\%$ ceiling
for the product of the input and output determinants.  For the balanced mixed-semigroup
determinant, the exact asymptotic min-plus divisor is governed by a two-dimensional Gini norm.
The resulting universal ceiling is $4/5$; it improves to $7/10$ in the cross-coprime branch,
while the structural places $2$ and $3$ alone contribute only $3/10$ there.  Finally, using
the complete triangular $(n,k)$ grid does not remove the termwise obstruction: optimal
transport gives the universal ceiling $\pi/4$ relative to the sharp largest-term
archimedean bound.

These results identify the remaining opportunity precisely.  A determinant proof must obtain
leading-order divisibility from genuine cancellation among equal-valuation terms, not merely
from entrywise valuations; or it must beat the standard archimedean bound through the
Harish--Chandra--Itzykson--Zuber free energy.  The report gives exact formalization targets for
both possibilities.
\end{minipage}
\vfill
{\large Prepared for Vladimir Reshetnikov\par}
\vspace{0.12cm}
{\normalsize OpenAI GPT-5.6 Pro\par}
\vspace{0.35cm}
{\large August 21, 2026\par}
\end{titlepage}

\pagenumbering{roman}

\begin{abstract}
Let $\vartheta=\log 3/\log 2$.  The Alaoglu--Erd\H{o}s conjecture says that
$2^x,3^x\in\Z$ for real $x$ forces $x\in\Z$.  The attached unified report develops an exact
conditional structure: if the conjecture fails, there is a least nonintegral generator
$\beta$, integers $M=2^\beta$ and $A=3^\beta$, and unique coordinate descriptions
$\Sol=\Nzero+\Nzero\beta$ and $\Cand=\{2^nM^k:n,k\in\Nzero\}$.  It also gives local and mixed
Vandermonde determinant bounds, but finds a missing logarithm between the number of nodes
available in a dyadic block and the number required by a purely archimedean interpolation
argument.

This continuation first normalizes $\beta=B+\alpha$, $0<\alpha<1$, and proves that failure is
equivalent to a single irrational lower mechanical word
$s_k=\lfloor(k+1)\alpha\rfloor-\lfloor k\alpha\rfloor$ simultaneously driving
\[
 R_{k+1}=\frac{(M/2^B)R_k}{2^{s_k}},\qquad
 S_{k+1}=\frac{(A/3^B)S_k}{3^{s_k}},
\]
with $R_k\in[1,2)$ and $S_k\in[1,3)$.  This is a synchronized dyadic--triadic Sturmian
formulation, and the primitive-output theorem makes at least one denominator sequence attain
its full base-power depth.  Next, Baker--W\"ustholz/Matveev lower bounds give effective
constants $c(M)>0$ and $\mu(M)<\infty$ for which
$|\beta-p/q|\ge c(M)q^{-\mu(M)}$; consequently the convergent denominators grow only
polynomially from one stage to the next.  This is much stronger than the elementary
single-exponential recurrence in the source report, although still compatible with the exact
conditional block population.

The main new work is a sharp $p$-adic audit.  For the ordinary Vandermonde on all conditional
candidates in one dyadic block, all valuations at primes supporting $M$ can be evaluated
exactly, and the analogous output formula is exact up to $O(T^2)$.  If
$a=v_2(M)$, $b=v_3(A)$, $\delta_2=\beta-a$, and $\delta_3=\beta-b$, then the all-support
termwise divisor of the product of the input and output Vandermondes occupies the asymptotic
fraction
\[
 1-\frac{\delta_2\log2+\delta_3\log3}{3\beta\log6}
\]
of the leading archimedean budget.  Since the primitive pair has $\min(a,b)=0$, this fraction
is at most $1-\log2/(3\log6)=0.871049\ldots$.

For the mixed-semigroup determinant, let $(u_j,v_j)$ be the first $N$ points of
$\Nzero^2$ ordered by $u+\vartheta v$.  The first-order assignment minima at every prime admit
an exact continuum limit.  If $X=(X_1,X_2)$ is uniform on the triangle
$\{x,z\ge0:x+z\le1\}$, define the Gini norm
$\mathfrak G(r,s)=\mathbb E|r(X_1-X_1')+s(X_2-X_2')|$.  Normalized prime-valuation slope
vectors $c_p\in\R^2$ satisfy $\sum_pc_p=0$, and the termwise divisor fraction is
\[
 1-\frac38\sum_p\mathfrak G(c_p)+o(1).
\]
The primitive condition and the triangle inequality give
$\sum_p\mathfrak G(c_p)\ge8/15$, hence the universal ceiling $4/5$.  If
$\gcd(MA,6)=1$, the ceiling improves to $7/10$; using only $p=2,3$ gives $3/10$.  The Gini norm
is evaluated in closed form from the three vertex values $0,r,s$.  This corrects a subsequent
repository note whose simultaneous base-depth normalization and proposed all-prime $2/3$
ceiling are not valid in full generality.

Finally, the complete triangular grid of all $(n,k)$ with $n+k\beta\le T$ supplies
$\asymp T^2$ rows.  Nevertheless, termwise divisibility still has a universal transport
ceiling.  The normalized row heights and mixed column weights both have quantile $Q(t)=\sqrt t$;
therefore the smallest and largest assignment energies tend to $\pi/8$ and $1/2$,
respectively.  Any all-prime min-plus divisor is bounded by the former, whereas the sharp
largest-term determinant bound is governed by the latter.  Their ratio is $\pi/4$.  The only
remaining determinant mechanism is thus genuine same-valuation cancellation or a sharper
archimedean free-energy estimate.  The Harish--Chandra--Itzykson--Zuber identity is recorded
as the natural analytic framework, and the report ends with a branch-sensitive research and
Lean formalization program.
\end{abstract}

\tableofcontents
\clearpage
\pagenumbering{arabic}

\section{Executive answer, scope, and trust boundaries}

\subsection{The requested goal and the outcome}

The requested primary goal was to continue the attached research program and prove the
following statement.

\begin{conjecture}[Alaoglu--Erd\H{o}s]
If $x\in\R$ and $2^x,3^x\in\Z$, then $x\in\Z$.
\end{conjecture}

I did not find an unconditional proof.  In particular, none of the arguments below converts
\[
  \log(2^x)\log3-\log(3^x)\log2=0
\]
into a forbidden linear form in logarithms of algebraic numbers.  The equality is bilinear in
logarithms, and normalizing it gives exactly the unresolved $2\times2$ determinant behind the
four exponentials conjecture.  Baker-type estimates control every \emph{nonzero} linear form
that occurs after fixing integer coefficients; they do not prove that this special bilinear
form cannot vanish.

The continuation nevertheless changes the technical map in useful ways.  It obtains an exact
symbolic-dynamical formulation, replaces an elementary continued-fraction estimate by an
effective polynomial one, computes the full leading $p$-adic content of ordinary block
Vandermondes, and gives sharp continuum formulas for the termwise divisibility of two much
larger determinant constructions.  These calculations do more than say that a method ``seems
too weak'': they quantify the leading constant that is available and isolate the precise
kind of extra cancellation a successful proof would have to produce.

\subsection{The principal new results}

\begin{longtable}{@{}>{\raggedright\arraybackslash}p{0.25\textwidth}>{\raggedright\arraybackslash}p{0.55\textwidth}>{\raggedright\arraybackslash}p{0.13\textwidth}@{}}
\toprule
Result & Content & Status\\
\midrule
\endfirsthead
\toprule
Result & Content & Status\\
\midrule
\endhead
Synchronized Sturmian orbit & A counterexample is equivalent to one irrational mechanical word simultaneously controlling the dyadic and triadic denominator jumps of two compact rational recurrences. & \statusnew\\
\addlinespace
Effective irrationality of $\beta$ & Baker--W\"ustholz/Matveev gives $|\beta-p/q|\ge c(M)q^{-\mu(M)}$ and hence polynomial, rather than exponential, recurrence bounds for continued-fraction denominators. & \statusclassical{} plus \statusnew\\
\addlinespace
Ordinary Vandermonde audit & The valuations at every prime in the fixed supports of $M$ and $A$ are computed to leading order.  The product construction has universal ceiling $1-\log2/(3\log6)=0.871049\ldots$. & \statusnew\\
\addlinespace
Gini assignment formula & For the balanced mixed-semigroup determinant, the exact termwise all-prime divisor fraction is $1-\frac38\sum_p\mathfrak G(c_p)+o(1)$ for an explicit norm $\mathfrak G$. & \statusnew\\
\addlinespace
Universal $4/5$ ceiling & The primitive condition forces $\sum_p\mathfrak G(c_p)\ge8/15$.  Thus entrywise/min-plus divisibility can supply at most $80\%$ of the leading block determinant height. & \statusnew\\
\addlinespace
Cross-coprime constants & If $\gcd(MA,6)=1$, the all-prime ceiling is $7/10$, and the places $2$ and $3$ alone contribute exactly $3/10$ in the continuum limit. & \statusnew\\
\addlinespace
Full-grid transport ceiling & On the complete triangular $(n,k)$ grid, any termwise all-prime divisor is bounded by the countermonotone energy $\pi/8$, while the largest determinant term has energy $1/2$; the ratio is $\pi/4$. & \statusnew\\
\addlinespace
HCIZ reformulation & The full-grid archimedean problem is expressed by the Harish--Chandra--Itzykson--Zuber integral, identifying a concrete free-energy inequality as the remaining analytic target. & \statusclassical{} plus \statusnew\\
\bottomrule
\end{longtable}

\subsection{Evidence labels}

The labels used in this document are intentionally conservative.

\begin{description}
\item[\statuslean] The attached report says the statement has been accepted by the Lean kernel in the repository snapshot it audits.  I did not independently replay the entire Lean build in this environment.
\item[\statusrepo] A current repository commit or report states that an aggregate build passed.  This is repository metadata, not an independently reproduced kernel audit here.
\item[\statusnew] A complete mathematical argument is supplied in this report.  It has been algebraically and numerically stress-tested, but it is not yet formalized or peer-reviewed.
\item[\statusclassical] The result is standard literature, used with an explicit citation.
\item[\statuscomp] A finite calculation or convergence table checks formulas and constants; it is not used as a substitute for proof.
\item[\statusopen] The statement is proposed as a precise target, not asserted.
\end{description}

\subsection{What the determinant ceilings do and do not prove}

A ceiling such as $4/5$ has a specific meaning.  Given an integer determinant
$D=\det(E_{ij})$, define at each prime $p$ the minimum valuation of a Leibniz term,
\[
  \tau_p=\min_{\sigma}\sum_i v_p(E_{i,\sigma(i)}).
\]
Every term, and therefore $D$, is divisible by $p^{\tau_p}$.  The product
$\prod_pp^{\tau_p}$ is the \emph{termwise} or \emph{tropical} divisor.  It ignores extra
valuation caused by cancellation among several terms attaining the same minimum.  The $4/5$
theorem says that this automatic divisor cannot occupy more than four-fifths of the leading
archimedean height for the balanced block construction.  It does \emph{not} say that the true
valuation $v_p(D)$ is at most the tropical minimum; the opposite inequality holds, and extra
cancellation may be large.  Thus the result retires only the naïve entrywise-divisibility
version of the method.  It leaves a sharply defined cancellation problem.

\section{Baseline inherited from the unified report}

\subsection{Notation}

Set
\[
  \vartheta=\frac{\log3}{\log2}.
\]
The solution and candidate sets are
\[
  \Sol=\{x\ge0:2^x,3^x\in\N\},
  \qquad
  \Cand=\{m\in\N:m^{\vartheta}\in\N\}.
\]
The map $x\mapsto2^x$ is a bijection $\Sol\to\Cand$.  Integer exponents give the obvious
solutions $\Nzero\subseteq\Sol$ and candidates $2^{\Nzero}\subseteq\Cand$.

Throughout the conditional sections, suppose the conjecture fails.  The exact structure
proved in the source report then supplies a least nonintegral solution $\beta>0$ and primitive
outputs
\[
  M=2^\beta\in\N,
  \qquad
  A=3^\beta\in\N.
\]
Every solution and candidate has a unique coordinate representation
\begin{equation}\label{eq:exact-structure}
  \Sol=\{n+k\beta:n,k\in\Nzero\},
  \qquad
  \Cand=\{2^nM^k:n,k\in\Nzero\}.
\end{equation}
The associated output is
\[
  (2^nM^k)^\vartheta=3^nA^k.
\]
We write
\[
  a=v_2(M),\qquad b=v_3(A),\qquad c=v_2(A),\qquad d=v_3(M).
\]
The inherited primitive-output theorem includes the crucial condition
\begin{equation}\label{eq:no-matched-depth}
  \min(a,b)=0.
\end{equation}
It is important not to silently strengthen \eqref{eq:no-matched-depth} to $a=b=0$.

\subsection{Inherited results used here}

The following facts are taken from the attached unified report and its referenced Lean
modules.

\begin{itemize}
\item \statuslean{} A nonintegral solution, if it exists, is transcendental.  The constant
      $\vartheta$ is transcendental, by Gelfond--Schneider.
\item \statuslean{} Failure gives the exact free rank-two monoid structure
      \eqref{eq:exact-structure}, including uniqueness of $(n,k)$.
\item \statuslean{} In a dyadic exponent block $[T,T+1)$, the solutions are exactly
      \[
        T
        \quad\text{and}\quad
        T-\lfloor k\beta\rfloor+k\beta,
        \qquad 1\le k\le\left\lfloor\frac{T+1}{\beta}\right\rfloor.
      \]
      Thus the block count is $1+\lfloor(T+1)/\beta\rfloor$.
\item \statuslean{} The primitive pair obeys \eqref{eq:no-matched-depth}; both outputs cannot
      simultaneously be $\{2,3\}$-smooth.
\item \statuslean{} Mixed monomials $m^u(m^\vartheta)^v$ are integers for
      $(u,v)\in\Nzero^2$, and generalized Vandermonde determinants built from distinct
      candidate nodes and distinct exponents $u+\vartheta v$ are nonzero integers.
\item \statuslean{} The kernel-checked finite exclusion in the attached report gives
      $x\ge12$ for a nonintegral solution.  The report separately records a directed-rounding
      software exclusion through $x\le26$.
\end{itemize}

The continuation does not reprove this large inherited corpus.  It restates every inherited
fact at the point where it enters a new argument.

\subsection{Repository reconciliation as of August 21, 2026}

The attached source is internally newer than one bibliographic pin printed near its end: its
body and status matrix already discuss the enlarged algebraic-output and determinant module
surface, although a repository item still names commit
\sha{5d3982ceb12644fdd3499569d9db5f62f9a31dc6}.  During this continuation I checked the
current public repository history.  Commit
\sha{78f80c5d3f6f9afd64ecb72ddb154c85e2036bc1} reports a successful aggregate build with
$67$ imported modules and merges the later algebraic-output work.  Subsequent commits add a
literature compass and a short $p$-adic priority note; the latest snapshot consulted while
writing was \sha{ff9a9f5aa3a3f02599de93fc053e60918f269855}.

\begin{statusbox}{\statusrepo}
The build claims in the preceding paragraph are taken from repository commit messages and
files.  I did not independently run all Lean jobs.  The mathematical arguments newly stated
below are therefore labelled \statusnew, not \statuslean.
\end{statusbox}

\subsection{The exact remaining logarithmic wall}

Writing $M=2^\beta$ and $A=3^\beta$ gives
\[
  \beta=\frac{\log M}{\log2}=\frac{\log A}{\log3}
\]
and hence
\begin{equation}\label{eq:bilinear-wall}
  \log M\,\log3-\log A\,\log2=0.
\end{equation}
The four numbers $\log2,\log3,\log M,\log A$ are logarithms of algebraic numbers.  Equation
\eqref{eq:bilinear-wall} is a vanishing $2\times2$ determinant.  Four exponentials would
exclude it because $2,3$ are multiplicatively independent and $1,\beta$ are linearly
independent over $\Q$, but four exponentials remains open.  The new work below therefore
pursues information that is invisible in the bare real logarithmic determinant: exact floor
coding, effective separation, and valuations of large integer determinants.

\section{A synchronized Sturmian model of a counterexample}
\label{sec:sturmian}

\subsection{Fractional-part normalization}

Write
\[
  \beta=B+\alpha,
  \qquad B=\lfloor\beta\rfloor\in\Nzero,
  \qquad 0<\alpha<1.
\]
Since a nonintegral solution is transcendental, $\alpha$ is irrational.  Define the rational
localized multipliers
\begin{equation}\label{eq:UV}
  U=\frac{M}{2^B}=2^\alpha\in\Q\cap(1,2),
  \qquad
  V=\frac{A}{3^B}=3^\alpha\in\Q\cap(1,3).
\end{equation}
For $k\ge0$, put
\[
  e_k=\lfloor k\alpha\rfloor,
  \qquad
  s_k=e_{k+1}-e_k\in\{0,1\}.
\]
The binary sequence $s=s_0s_1s_2\cdots$ is the lower mechanical word of slope $\alpha$.
Define compact normalized orbits
\begin{equation}\label{eq:normalized-orbits}
  R_k=\frac{U^k}{2^{e_k}}=2^{\fract{k\alpha}}\in[1,2),
  \qquad
  S_k=\frac{V^k}{3^{e_k}}=3^{\fract{k\alpha}}\in[1,3).
\end{equation}
Both orbits use the \emph{same} fractional part and hence the same symbolic itinerary.

\begin{theorem}[Synchronized dyadic--triadic Sturmian recurrence]\label{thm:sturmian-recurrence}
Assume a primitive nonintegral solution $\beta$ exists.  Then the rational sequences
$R_k,S_k$ in \eqref{eq:normalized-orbits} satisfy
\begin{equation}\label{eq:sturmian-recurrence}
  R_0=S_0=1,
  \qquad
  R_{k+1}=\frac{U R_k}{2^{s_k}},
  \qquad
  S_{k+1}=\frac{V S_k}{3^{s_k}},
\end{equation}
where
\begin{equation}\label{eq:sturmian-threshold}
  s_k=1
  \quad\Longleftrightarrow\quad
  \fract{k\alpha}\ge1-\alpha
  \quad\Longleftrightarrow\quad
  UR_k\ge2
  \quad\Longleftrightarrow\quad
  VS_k\ge3.
\end{equation}
Moreover,
\begin{equation}\label{eq:prefix-count}
  \sum_{j=0}^{N-1}s_j=\lfloor N\alpha\rfloor
  \qquad(N\ge1).
\end{equation}
\end{theorem}

\begin{proof}
The recurrence follows by subtracting consecutive floor exponents:
\[
  R_{k+1}
  =\frac{U^{k+1}}{2^{e_{k+1}}}
  =\frac{U U^k}{2^{e_k+s_k}}
  =\frac{UR_k}{2^{s_k}},
\]
and the same calculation with base $3$ gives the formula for $S_k$.  Since
$e_{k+1}=e_k+1$ exactly when $\fract{k\alpha}+\alpha\ge1$, the first equivalence in
\eqref{eq:sturmian-threshold} holds.  The identities
$UR_k=2^{\alpha+\fract{k\alpha}}$ and
$VS_k=3^{\alpha+\fract{k\alpha}}$ give the other two.  Finally,
\eqref{eq:prefix-count} telescopes:
$\sum_{j<N}s_j=e_N-e_0=\lfloor N\alpha\rfloor$.
\end{proof}

\subsection{Sturmian consequences}

The following facts are classical consequences of irrational mechanical coding
\cite{MorseHedlund1940,Lothaire2002}.

\begin{corollary}[Aperiodicity, balance, and minimal complexity]\label{cor:sturmian}
The word $s$ in Theorem~\ref{thm:sturmian-recurrence} is aperiodic.  Any two factors of equal
length contain numbers of $1$'s differing by at most one, and the number of distinct factors
of length $n$ is exactly $n+1$.
\end{corollary}

Aperiodicity has an elementary proof in this setting.  If $s$ were eventually periodic, then
its limiting density of $1$'s would be rational.  Equation \eqref{eq:prefix-count} shows that
this density is $\alpha$, contradicting irrationality.  Balance and factor complexity are the
standard Sturmian characterization.

The point is not merely terminological.  A hypothetical counterexample is no longer described
only by two rational numbers $U,V$ satisfying a logarithmic equality.  It determines one
minimal-complexity aperiodic word, and that same word must be realized by two rational
piecewise-multiplicative systems with multiplicatively independent reset bases $2$ and $3$.
This is a much more rigid synchronized object.

\subsection{Exact reduced denominators}

The primitive factor depths determine how much cancellation occurs in
\eqref{eq:normalized-orbits}.  Write
\[
  M=2^aW,\quad W\text{ odd},
  \qquad
  A=3^bZ,\quad 3\nmid Z.
\]
Because $M$ is not a power of $2$ and $A$ is not a power of $3$, one has $W,Z>1$.

\begin{proposition}[Denominator exponents and common jump word]\label{prop:denominators}
In lowest terms,
\begin{align}
  R_k&=\frac{W^k}{2^{q^{(2)}_k}},
  &q^{(2)}_k&=\lfloor k\beta\rfloor-ak
              =k(B-a)+\lfloor k\alpha\rfloor,\label{eq:q2}\\
  S_k&=\frac{Z^k}{3^{q^{(3)}_k}},
  &q^{(3)}_k&=\lfloor k\beta\rfloor-bk
              =k(B-b)+\lfloor k\alpha\rfloor.\label{eq:q3}
\end{align}
Consequently
\[
  q^{(2)}_{k+1}-q^{(2)}_k=B-a+s_k,
  \qquad
  q^{(3)}_{k+1}-q^{(3)}_k=B-b+s_k.
\]
At least one of the two sequences has baseline jump $B$, because $\min(a,b)=0$.
\end{proposition}

\begin{proof}
Using $\lfloor k\beta\rfloor=kB+\lfloor k\alpha\rfloor$,
\[
 R_k=\frac{M^k}{2^{\lfloor k\beta\rfloor}}
     =\frac{2^{ak}W^k}{2^{\lfloor k\beta\rfloor}}
     =\frac{W^k}{2^{\lfloor k\beta\rfloor-ak}}.
\]
The numerator is odd, so this is reduced.  The base-$3$ calculation is identical.  Taking
successive differences gives the jump formulas, and the last claim is the inherited
primitive condition \eqref{eq:no-matched-depth}.
\end{proof}

Thus the same Sturmian bit $s_k$ records whether \emph{both} reduced denominators take their
large or small permitted jump at time $k$.  One denominator can have a shallower constant
baseline, but the aperiodic fluctuation is synchronized exactly.

\subsection{An exact equivalent exclusion principle}

The source report already contains localization reformulations such as
$2^\alpha\in\Z[1/2]$ and $3^\alpha\in\Z[1/3]$.  The recurrence packages them into a symbolic
statement.

\begin{theorem}[Synchronized Sturmian exclusion is equivalent to Alaoglu--Erd\H{o}s]
\label{thm:sturmian-equivalence}
The Alaoglu--Erd\H{o}s conjecture is equivalent to the nonexistence of data
\[
  \alpha\in(0,1)\setminus\Q,
  \qquad
  U\in\Q\cap(1,2),
  \qquad
  V\in\Q\cap(1,3)
\]
such that, for the lower mechanical word
$s_k=\lfloor(k+1)\alpha\rfloor-\lfloor k\alpha\rfloor$, the recurrences
\[
  R_0=S_0=1,
  \qquad R_{k+1}=UR_k/2^{s_k},
  \qquad S_{k+1}=VS_k/3^{s_k}
\]
remain in $[1,2)\cap\Q$ and $[1,3)\cap\Q$, respectively, for every $k$, with
$U=2^\alpha$ and $V=3^\alpha$.
Equivalently, it suffices to rule out irrational $\alpha$ for which
\[
  2^\alpha\in\Z[1/2],
  \qquad
  3^\alpha\in\Z[1/3].
\]
\end{theorem}

\begin{proof}
A counterexample $\beta=B+\alpha$ gives the data by
Theorem~\ref{thm:sturmian-recurrence}.  Conversely, if such $\alpha,U,V$ exist, choose an
integer $B$ at least as large as the powers of $2$ and $3$ occurring in the respective
reduced denominators of $U$ and $V$.  Then
\[
  2^{B+\alpha}=2^BU\in\N,
  \qquad
  3^{B+\alpha}=3^BV\in\N,
\]
while $B+\alpha\notin\Z$.  The recurrence identities are equivalent to
$R_k=2^{\fract{k\alpha}}$ and $S_k=3^{\fract{k\alpha}}$, so they add an exact dynamic coding
but no hidden hypothesis.
\end{proof}

\subsection{Cobham-type and Mahler-type diagnostics}

Cobham's theorem says that a sequence automatic in two multiplicatively independent integer
bases is ultimately periodic \cite{Cobham1972,AlloucheShallit2003}.  Sturmian words of
irrational slope are not automatic in any integer base.  Theorem~\ref{thm:sturmian-equivalence}
therefore suggests a tempting strategy: show that the dyadic rational recurrence makes $s$
$2$-automatic and that the triadic recurrence makes it $3$-automatic.

That implication is \emph{not} presently available.  Rationality of the multiplier does not
turn the compact rational state $R_k$ or $S_k$ into a finite-state object; the exact reduced
denominators grow without bound.  Cobham's theorem becomes applicable only after proving a
finite-kernel or finite-substitution property, and that is essentially a new arithmetic
statement.  The synchronized formulation is useful because it identifies this missing lemma
cleanly rather than treating ``automaticity'' as a metaphor.

\begin{target}[Finite-kernel extraction]
Show that simultaneous rational realizability of the same irrational mechanical word by
\eqref{eq:sturmian-recurrence} forces a finite $2$-kernel and a finite $3$-kernel, or derive a
weaker pair of incompatible substitutive structures.
\end{target}

A Mahler-function variant is also conceivable: encode the jump word by
$F(z)=\sum_{k\ge0}s_kz^k$.  Aperiodicity and the finite coefficient set already imply, by
classical results of Fatou type, that $F$ is transcendental over $\Q(z)$.  A proof would have
to derive incompatible functional equations from the two rational recurrences.  No such
equations were found here; the recurrence is multiplicative in the state, not self-similar in
the index.

\section{Effective logarithmic separation}
\label{sec:baker}

\subsection{From linear forms in two logarithms to an irrationality measure}

The source report proves an elementary denominator-growth statement: a very good rational
approximation to a hypothetical solution $\beta$ would force a correspondingly large output,
and consecutive continued-fraction denominators obey a single-exponential recurrence.  Since
$\beta$ is itself a ratio of logarithms of fixed integers, Baker's theory gives a substantially
stronger conclusion.

\begin{theorem}[Effective finite irrationality measure for the primitive generator]
\label{thm:baker-beta}
Assume a primitive counterexample $\beta$ exists, and put $M=2^\beta$.  There are effectively
computable constants $c_M>0$ and $\mu_M<\infty$ such that, for every $p\in\Z$ and $q\ge2$,
\begin{equation}\label{eq:beta-irrat-measure}
  \left|\beta-\frac pq\right|\ge c_Mq^{-\mu_M}.
\end{equation}
One may choose $\mu_M$ bounded by an effectively computable absolute constant times
$1+\log M$, hence by an absolute constant times $1+\beta$.  The same conclusion follows from
$\beta=\log A/\log3$.
\end{theorem}

\begin{proof}
Consider the linear form
\[
  \Lambda=q\log M-p\log2.
\]
It is nonzero: otherwise $M^q=2^p$, which would make $\beta=p/q$ rational; a rational solution
is integral by the elementary rational-exponent classification already formalized in the
source corpus.  The Baker--W\"ustholz theorem, or Matveev's explicit lower bound for a
homogeneous linear form in logarithms of algebraic numbers, gives
\[
  \log|\Lambda|\ge-C\,(1+\log M)\log H-C'(1+\log M),
\]
where $H\ge\max(3,|p|,q)$ and $C,C'$ are effective absolute constants after the fixed number
$2$ is absorbed.  For approximations with $|p|\le(\beta+1)q$, one has $H\ll_Mq$; for all other
$p$ the desired inequality is trivial after decreasing $c_M$.  Since
\[
  \left|\beta-\frac pq\right|=\frac{|\Lambda|}{q\log2},
\]
absorbing fixed factors and one additional power of $q$ gives
\eqref{eq:beta-irrat-measure}.  Matveev's dependence on the logarithmic heights is linear in
$\log M=\beta\log2$ here, which gives the final scale assertion.
\end{proof}

\begin{statusbox}{\statusclassical{} plus \statusnew}
The lower-bound theorem for logarithmic forms is classical
\cite{Baker1975,BakerWustholz1993,Matveev2000}.  Its application to the primitive generator,
the continued-fraction consequences below, and the comparison with the source report are
new deductions in this continuation.  No useful small numerical value of $\mu_M$ is claimed;
standard explicit constants are very large.
\end{statusbox}

\subsection{Polynomial recurrence for convergent denominators}

Let $p_n/q_n$ be the continued-fraction convergents of $\beta$.  The standard upper bound
\[
  \left|\beta-\frac{p_n}{q_n}\right|<\frac1{q_nq_{n+1}}
\]
combined with \eqref{eq:beta-irrat-measure} gives the following.

\begin{corollary}[Polynomial denominator growth]\label{cor:q-growth}
There is an effective constant $C_M>0$ such that
\begin{equation}\label{eq:q-growth}
  q_{n+1}\le C_M q_n^{\mu_M-1}
\end{equation}
for every $n$ with $q_n\ge2$.  Consequently the next partial quotient satisfies
$a_{n+1}=O_M(q_n^{\mu_M-2})$.
\end{corollary}

This is qualitatively stronger than the inherited elementary estimate
$q_{n+1}<2M^{q_n}$.  It rules out Liouville behavior and says that the exceptionally long
near-repetitions of the Sturmian word are polynomially controlled in the preceding standard
word length.  The exponent depends on the unknown integer $M$, however, and may be enormous.

\subsection{Polynomial spacing in conditional blocks}

The irrationality measure also upgrades the separation of the block nodes.  From
\eqref{eq:beta-irrat-measure}, for $d\ge1$ and the nearest integer $r$ to $d\beta$,
\begin{equation}\label{eq:beta-mod-one}
  \distZ{d\beta}=|d\beta-r|
  \ge c_M d^{1-\mu_M}.
\end{equation}

For an integer $T\ge0$, put
\[
  K_T=\left\lfloor\frac{T+1}{\beta}\right\rfloor,
  \qquad
  x_{T,k}=T-\lfloor k\beta\rfloor+k\beta
         =T+\fract{k\beta},
  \quad 0\le k\le K_T.
\]
The associated candidate and output nodes are
\begin{equation}\label{eq:block-pairs}
  m_{T,k}=2^{x_{T,k}}=2^{T-\lfloor k\beta\rfloor}M^k,
  \qquad
  a_{T,k}=3^{x_{T,k}}=3^{T-\lfloor k\beta\rfloor}A^k.
\end{equation}

\begin{corollary}[Effective block separation]\label{cor:block-separation}
For distinct $i,j\in\{0,\ldots,K_T\}$,
\begin{align}
 |x_{T,j}-x_{T,i}|&\ge c_MK_T^{1-\mu_M},\label{eq:x-sep}\\
 |m_{T,j}-m_{T,i}|&\ge (\log2)c_M2^T K_T^{1-\mu_M},\label{eq:m-sep}\\
 |a_{T,j}-a_{T,i}|&\ge (\log3)c_M3^T K_T^{1-\mu_M}.
 \label{eq:a-sep}
\end{align}
\end{corollary}

\begin{proof}
The ordinary distance between two fractional parts is at least their circular distance, so
\[
 |\fract{j\beta}-\fract{i\beta}|
 \ge\distZ{(j-i)\beta}
 \ge c_M|j-i|^{1-\mu_M}
 \ge c_MK_T^{1-\mu_M},
\]
because $1-\mu_M<0$.  This proves \eqref{eq:x-sep}.  The mean value theorem applied to $2^x$
and $3^x$ on $[T,T+1]$ gives the remaining inequalities.
\end{proof}

\subsection{Separation on the full \texorpdfstring{$(n,k)$}{(n,k)} lattice}

The same estimate applies to the two-dimensional conditional exponent grid.

\begin{corollary}[Polynomial separation of conditional exponents]\label{cor:grid-separation}
There are effective $c_M,\mu_M>0$ such that, for distinct integer pairs
$(n,k),(n',k')$ with $0\le k,k'\le H$,
\[
  |(n+k\beta)-(n'+k'\beta)|
  \ge \min\{1,c_MH^{1-\mu_M}\}.
\]
\end{corollary}

\begin{proof}
If $k=k'$, the difference is a nonzero integer.  Otherwise divide the linear form
$(n-n')+(k-k')\beta$ by $|k-k'|$ and apply
\eqref{eq:beta-irrat-measure}.
\end{proof}

The column exponent lattice $u+\vartheta v$ has an analogous effective polynomial spacing
with absolute constants, because $\vartheta=\log3/\log2$ is a fixed ratio of two logarithms.
More refined explicit estimates are known for linear forms in $\log2$ and $\log3$; for
example Wu's simultaneous linear-independence measure leads, in the usual normalization, to
an irrationality exponent for $\vartheta$ below approximately $8.62$
\cite{Wu2003}.  The exact convention-dependent numerical conversion is not needed below.

\subsection{Why the Baker upgrade does not close the conjecture}

The effective spacing is useful, but its scale is still compatible with the exact conditional
population.  A dyadic exponent block contains only $K_T+1\asymp T$ nodes.  A polynomial lower
bound on their minimum spacing imposes no contradiction on $T$ points in a fixed interval.
Likewise, the full triangular grid has $\asymp T^2$ points in an interval of exponent length
$T$, and polynomially small gaps are expected.

Baker's theorem therefore strengthens compactness and determinant estimates but does not
supply the missing logarithm by itself.  Its most promising roles are indirect:

\begin{itemize}
\item controlling the smallest row and column gaps in an exact Harish--Chandra--Itzykson--Zuber
      factorization;
\item bounding how often a $p$-adic assignment problem can be nearly degenerate;
\item limiting the lengths of near-periodic Sturmian blocks that a congruence argument must
      handle;
\item making computational searches over continued-fraction branches provably finite once a
      candidate output $M$ is fixed.
\end{itemize}

\section{Exact all-support valuations of ordinary block Vandermondes}

\subsection{The block determinants}

Fix an integer $T$ and abbreviate $K=K_T$.  Let $m_k=m_{T,k}$ and $a_k=a_{T,k}$ be the
$K+1$ input--output pairs in \eqref{eq:block-pairs}.  Their ordinary Vandermonde determinants
are
\begin{align}
  V_T^{(2)}&=\det(m_i^j)_{0\le i,j\le K}
            =\prod_{0\le i<j\le K}(m_j-m_i),\label{eq:Vm}\\
  V_T^{(3)}&=\det(a_i^j)_{0\le i,j\le K}
            =\prod_{0\le i<j\le K}(a_j-a_i).
  \label{eq:Va}
\end{align}
They are nonzero integers.  Since $m_k\in[2^T,2^{T+1})$ and
$a_k\in[3^T,3^{T+1})$,
\begin{align}
  \log|V_T^{(2)}|&\le R_TT\log2,\label{eq:Vm-arch}\\
  \log|V_T^{(3)}|&\le R_T(T\log3+\log2),\label{eq:Va-arch}
\end{align}
where
\[
  R_T=\binom{K+1}{2}=\frac{K(K+1)}2.
\]
The additive $R_T\log2$ in \eqref{eq:Va-arch} is lower order.

\subsection{The input determinant: an exact formula at every support prime}

Factor
\[
  M=2^aW,
  \qquad W\text{ odd},
  \qquad
  \delta_2=\beta-a=\log_2W.
\]
As noted above, $W>1$, so $W\ge3$ and
\begin{equation}\label{eq:delta2-gt1}
  \delta_2\ge\log_2 3>1.
\end{equation}
Because $a$ is an integer,
\begin{equation}\label{eq:m-delta-form}
  m_k=2^{T-\lfloor k\delta_2\rfloor}W^k.
\end{equation}
The $2$-adic valuations in \eqref{eq:m-delta-form} strictly decrease with $k$, while every
odd support-prime valuation strictly increases.

\begin{proposition}[Exact support valuations for $V_T^{(2)}$]
\label{prop:input-vander-valuations}
For $0\le i<j\le K$,
\begin{align}
  v_2(m_j-m_i)&=T-\lfloor j\delta_2\rfloor,
  \label{eq:v2-gap}\\
  v_p(m_j-m_i)&=i\,v_p(W)
  \qquad(p\mid W).
  \label{eq:vp-gap}
\end{align}
Consequently
\begin{align}
 v_2(V_T^{(2)})
 &=\sum_{j=1}^Kj\bigl(T-\lfloor j\delta_2\rfloor\bigr),
 \label{eq:v2-V}\\
 v_p(V_T^{(2)})
 &=v_p(W)\sum_{i=0}^{K-1}i(K-i)
 \qquad(p\mid W).
 \label{eq:vp-V}
\end{align}
\end{proposition}

\begin{proof}
Put $\rho_k=T-\lfloor k\delta_2\rfloor$.  By
\eqref{eq:delta2-gt1}, $\rho_i>\rho_j$ when $i<j$.  From
\eqref{eq:m-delta-form},
\[
 m_j-m_i
 =2^{\rho_j}W^i\left(W^{j-i}-2^{\rho_i-\rho_j}\right).
\]
The parenthesis is odd, proving \eqref{eq:v2-gap}.  For an odd prime $p\mid W$, the two terms
$m_i$ and $m_j$ have distinct valuations $i v_p(W)<jv_p(W)$, so the ultrametric equality gives
\eqref{eq:vp-gap}.  Summing over pairs gives \eqref{eq:v2-V}--\eqref{eq:vp-V}.
\end{proof}

Let $Q_T^{(2)}$ be the part of $V_T^{(2)}$ supported on the primes dividing $2M$.  Then
Proposition~\ref{prop:input-vander-valuations} gives the exact identity
\begin{equation}\label{eq:logQm-exact}
 \log Q_T^{(2)}
 =\log2\sum_{j=1}^Kj\bigl(T-\lfloor j\delta_2\rfloor\bigr)
  +\log W\sum_{i=0}^{K-1}i(K-i).
\end{equation}
The two elementary sums have asymptotics
\begin{align*}
 \sum_{j=1}^Kj\bigl(T-\lfloor j\delta_2\rfloor\bigr)
 &=\frac{TK^2}{2}-\frac{\delta_2K^3}{3}+O_\beta(T^2),\\
 \sum_{i=0}^{K-1}i(K-i)&=\frac{K^3}{6}+O(K^2).
\end{align*}
Using $\log W=\delta_2\log2$ and $K/T\to1/\beta$ yields
\begin{equation}\label{eq:Qm-ratio}
 \frac{\log Q_T^{(2)}}{R_TT\log2}
 \longrightarrow
 1-\frac{\delta_2}{3\beta}.
\end{equation}
This includes every prime appearing in the entries.  Any further prime divisor of the
Vandermonde comes from subtraction and hence from genuine cancellation beyond termwise
entry valuations.

\subsection{The output determinant and bounded tie corrections}

Factor
\[
  A=3^bZ,
  \qquad 3\nmid Z,
  \qquad
  \delta_3=\beta-b=\log_3Z>0.
\]
Then
\begin{equation}\label{eq:a-delta-form}
  a_k=3^{T-\lfloor k\delta_3\rfloor}Z^k.
\end{equation}
Unlike $\delta_2$, the number $\delta_3$ can be smaller than $1$; for example the least
possible base-free factor is $2$.  Equal consecutive $3$-adic valuations can therefore occur.
They contribute only a lower-order correction.

\begin{proposition}[Output support valuations]\label{prop:output-vander-valuations}
For every prime $p\mid Z$ and $0\le i<j\le K$,
\[
  v_p(a_j-a_i)=i\,v_p(Z).
\]
At $p=3$,
\[
 v_3(a_j-a_i)=T-\lfloor j\delta_3\rfloor
\]
after adding a nonnegative correction that is zero whenever
$\lfloor i\delta_3\rfloor<\lfloor j\delta_3\rfloor$.  The total correction over all pairs is
$O_{A,\beta}(K^2)$.  Consequently, if $Q_T^{(3)}$ denotes the part of
$V_T^{(3)}$ supported on primes dividing $3A$, then
\begin{equation}\label{eq:Qa-ratio}
  \frac{\log Q_T^{(3)}}{R_TT\log3}
  \longrightarrow
  1-\frac{\delta_3}{3\beta}.
\end{equation}
\end{proposition}

\begin{proof}
The assertion for $p\mid Z$ is the same distinct-valuation argument as before.  Put
$\rho_k=T-\lfloor k\delta_3\rfloor$.  If $\rho_i>\rho_j$, then
\[
 a_j-a_i=3^{\rho_j}Z^i
 \left(Z^{j-i}-3^{\rho_i-\rho_j}\right)
\]
and the parenthesis is not divisible by $3$, so the valuation is exactly $\rho_j$.  If
$\rho_i=\rho_j$, then
\[
 v_3(a_j-a_i)=\rho_j+v_3(Z^{j-i}-1).
\]
Equality of the floors implies $(j-i)\delta_3<1$, so $j-i$ belongs to a fixed finite set
depending only on $\delta_3$.  Hence the extra valuation is bounded by a constant depending on
$A$ and $\beta$.  There are $O(K^2)$ pairs.  The asymptotic calculation is then identical to
\eqref{eq:logQm-exact}.
\end{proof}

\subsection{The depth-sensitive product ceiling}

Combining \eqref{eq:Vm-arch}, \eqref{eq:Va-arch},
\eqref{eq:Qm-ratio}, and \eqref{eq:Qa-ratio} gives the principal conclusion of the ordinary
Vandermonde audit.

\begin{theorem}[All-support ordinary Vandermonde fraction]\label{thm:ordinary-fraction}
For a hypothetical primitive counterexample,
\begin{equation}\label{eq:ordinary-fraction}
 \frac{\log(Q_T^{(2)}Q_T^{(3)})}{R_TT\log6}
 \longrightarrow
 1-\frac{\delta_2\log2+\delta_3\log3}{3\beta\log6},
 \qquad
 \delta_2=\beta-a,
 \quad
 \delta_3=\beta-b.
\end{equation}
Because $\min(a,b)=0$, the limit is at most
\begin{equation}\label{eq:ordinary-universal-ceiling}
 1-\frac{\log2}{3\log6}
 =0.871049\,064\ldots.
\end{equation}
If $b=0$, the sharper ceiling is
\[
 1-\frac{\log3}{3\log6}=0.795599\,921\ldots.
\]
If $a=b=0$, the limit is exactly $2/3$.
\end{theorem}

\begin{proof}
Only the numerical ceiling remains to be shown.  If $a=0$, then $\delta_2=\beta$, so the
deficit in \eqref{eq:ordinary-fraction} is at least $\log2/(3\log6)$.  If $b=0$, it is at
least $\log3/(3\log6)$.  Taking the larger possible fraction gives
\eqref{eq:ordinary-universal-ceiling}.  If both vanish, the numerator of the deficit is
$\beta\log6$.
\end{proof}

\begin{table}[ht]
\centering
\caption{Leading termwise-divisor ceilings for ordinary block Vandermondes.}
\begin{tabular}{@{}lll@{}}
\toprule
Branch information & Exact/upper fraction & Decimal\\
\midrule
$a=b=0$ & $2/3$ & $0.666666\ldots$\\
$b=0$ & $\le1-\log3/(3\log6)$ & $0.795599\ldots$\\
No matched depth only & $\le1-\log2/(3\log6)$ & $0.871049\ldots$\\
\bottomrule
\end{tabular}
\end{table}

\subsection{Why even \texorpdfstring{$87\%$}{87 percent} is not enough}

The deficit in Theorem~\ref{thm:ordinary-fraction} is a fixed positive fraction once the
hypothetical primitive pair is fixed.  The archimedean height is $\asymp T^3$, whereas all
factorial, spacing, and tie corrections are $O_\beta(T^2\log T)$.  Thus no lower-order
refinement of the ordinary Vandermonde can close the remaining gap.

This conclusion is stronger than the earlier observation that the $2$-adic valuation alone
captures one third in the base-free case.  It includes every prime already present in the
entries and every possible base depth.  A successful ordinary-Vandermonde proof would need a
new leading-order source of divisibility in the prime factors of the \emph{differences}; the
fixed-support valuations are completely accounted for by
Theorem~\ref{thm:ordinary-fraction}.

\section{The mixed-semigroup determinant and its Gini norm}
\label{sec:mixed}

\subsection{Balanced low-weight columns}

Let
\[
  \Lambda=\{u+\vartheta v:(u,v)\in\Nzero^2\}.
\]
Because $\vartheta$ is irrational, every element has a unique pair $(u,v)$.  List the pairs
in increasing order of $u+\vartheta v$:
\[
  (u_0,v_0),(u_1,v_1),\ldots.
\]
For $N\ge1$, let
\[
  L_N=\max_{0\le j<N}(u_j+\vartheta v_j).
\]
The elementary lattice count
\begin{equation}\label{eq:semigroup-count}
 \#\{(u,v)\in\Nzero^2:u+\vartheta v\le L\}
 =\frac{L^2}{2\vartheta}+O_\vartheta(L)
\end{equation}
shows that
\begin{equation}\label{eq:L-asymp}
  L_N\sim\sqrt{2\vartheta N}.
\end{equation}
Moreover, the empirical distribution of
\begin{equation}\label{eq:scaled-cols}
  X_j=\left(\frac{u_j}{L_N},\frac{\vartheta v_j}{L_N}\right)
\end{equation}
converges to uniform area measure on the standard triangle
\begin{equation}\label{eq:triangle}
  \DeltaTri=\{(x,z)\in\R^2:x\ge0,
  z\ge0,
  x+z\le1\}.
\end{equation}
Here and below ``uniform'' means probability measure with density $2$ on this triangle.

For the block at height $T$, set $K=K_T$ and $N=K+1$.  Define
\begin{equation}\label{eq:mixed-D}
  D_T=\det\left(m_{T,k}^{u_j}a_{T,k}^{v_j}\right)_{0\le k,j<N}.
\end{equation}
The generalized-Vandermonde theorem inherited from the source report implies that $D_T$ is a
nonzero integer.  Put
\[
  B_j=2^{u_j}3^{v_j}.
\]
Since $m_{T,k}=2^{x_{T,k}}$ and $a_{T,k}=3^{x_{T,k}}$,
\begin{equation}\label{eq:entry-B}
  m_{T,k}^{u_j}a_{T,k}^{v_j}=B_j^{x_{T,k}}
  =B_j^{T+\fract{k\beta}}.
\end{equation}
Thus the Leibniz expansion gives the clean archimedean bound
\begin{equation}\label{eq:mixed-arch}
  \log|D_T|
  \le (T+1)\sum_{j<N}\log B_j+\log N!.
\end{equation}
The leading height is
\begin{equation}\label{eq:HT}
  H_T=T\sum_{j<N}\log B_j.
\end{equation}
By \eqref{eq:L-asymp} and the mean
$\mathbb E(X_1+X_2)=2/3$ on $\DeltaTri$,
\begin{equation}\label{eq:H-asymp}
 H_T\sim \frac23\,TNL_N\log2
 \asymp_\beta T^{5/2}.
\end{equation}
Both $\sum_j\log B_j$ and $\log N!$ in the error of \eqref{eq:mixed-arch} are
$o(H_T)$.

\subsection{The all-prime tropical divisor}

For a prime $p$, define the assignment minimum
\begin{equation}\label{eq:taup}
 \tau_{p,T}
 =\min_{\sigma\in S_N}
   \sum_{k=0}^{N-1}
   v_p\!\left(m_{T,k}^{u_{\sigma(k)}}a_{T,k}^{v_{\sigma(k)}}\right).
\end{equation}
Every Leibniz term is divisible by $p^{\tau_{p,T}}$, so
\begin{equation}\label{eq:Qtrop}
  Q_T^{\mathrm{trop}}=\prod_p p^{\tau_{p,T}}
  \quad\text{divides}\quad D_T.
\end{equation}
Only primes dividing $6MA$ occur in the product.

The problem is now to compute the leading asymptotic of
$\log Q_T^{\mathrm{trop}}/H_T$.  A crude estimate ``minimum $\le$ average'' loses the precise
branch dependence.  The exact answer is controlled by a simple probability norm.

\subsection{Antimonotone assignment and the Gini norm}

Let $X=(X_1,X_2)$ be uniform on $\DeltaTri$, and let $X'$ be an independent copy.  For
$c=(r,s)\in\R^2$, define
\begin{equation}\label{eq:Gini}
  \mathfrak G(r,s)
  =\mathbb E\left|r(X_1-X_1')+s(X_2-X_2')\right|.
\end{equation}
This is a norm on $\R^2$: homogeneity and the triangle inequality are immediate, and the
support of $X-X'$ spans $\R^2$, so only $(r,s)=(0,0)$ has norm zero.

Let $Y=rX_1+sX_2$, with ascending quantile $Q_Y(t)$.  Pairing an increasing uniform row
parameter with the decreasing ordering of $Y$ gives the continuum assignment functional
\begin{equation}\label{eq:Phi-def}
  \Phi(r,s)=\int_0^1 t\,Q_Y(1-t)\,dt.
\end{equation}

\begin{lemma}[Gini formula for the minimum assignment]\label{lem:gini-assignment}
For every $(r,s)\in\R^2$,
\begin{equation}\label{eq:Phi-G}
  \Phi(r,s)
  =\frac12\,\mathbb E Y-\frac14\,\mathfrak G(r,s)
  =\frac{r+s}{6}-\frac14\,\mathfrak G(r,s).
\end{equation}
\end{lemma}

\begin{proof}
The rearrangement inequality says that the minimum scalar product pairs the increasing row
parameter with the decreasing list of column values.  Changing variables $u=1-t$ in
\eqref{eq:Phi-def} gives
\[
  \Phi=\int_0^1(1-u)Q_Y(u)\,du.
\]
If $Y'$ is an independent copy, the minimum of $Y,Y'$ has quantile density
$2(1-u)$, so
\[
  \Phi=\frac12\mathbb E\min(Y,Y').
\]
Since
$\min(Y,Y')=(Y+Y'-|Y-Y'|)/2$, this becomes
$\frac12\mathbb EY-\frac14\mathbb E|Y-Y'|$.  Finally
$\mathbb EX_1=\mathbb EX_2=1/3$.
\end{proof}

\subsection{Closed form of the Gini norm}

The norm \eqref{eq:Gini} is elementary.  Sort the three vertex values
\[
  0,
  \qquad r,
  \qquad s
\]
as $\xi_0\le\xi_1\le\xi_2$, and put
\[
  A_0=\xi_1-\xi_0,
  \qquad
  B_0=\xi_2-\xi_1.
\]

\begin{proposition}[Explicit triangular Gini norm]\label{prop:G-explicit}
If $A_0+B_0>0$, then
\begin{equation}\label{eq:G-explicit}
 \mathfrak G(r,s)
 =\frac{2\left(2A_0^2+3A_0B_0+2B_0^2\right)}
        {15(A_0+B_0)}.
\end{equation}
At $(r,s)=(0,0)$ the value is $0$.  In particular,
\begin{equation}\label{eq:G-special}
  \mathfrak G(1,0)=\mathfrak G(0,1)
  =\mathfrak G(1,1)=\frac4{15},
  \qquad
  \mathfrak G(1,-1)=\frac7{15}.
\end{equation}
\end{proposition}

\begin{proof}
Translation of the three vertex values does not change a difference $Y-Y'$, so take
$\xi_0=0$, $\xi_1=A_0$, and $\xi_2=A_0+B_0=L$.  Projection of uniform area measure on a
triangle by a linear function has density
\[
 f(y)=
 \begin{cases}
 \displaystyle\frac{2y}{A_0L},&0\le y\le A_0,\\[6pt]
 \displaystyle\frac{2(L-y)}{B_0L},&A_0\le y\le L.
 \end{cases}
\]
Thus its distribution function is
$F(y)=y^2/(A_0L)$ on the first interval and
$F(y)=1-(L-y)^2/(B_0L)$ on the second.  The standard identity
\[
  \mathbb E|Y-Y'|=2\int_{-\infty}^{\infty}F(y)(1-F(y))\,dy
\]
gives \eqref{eq:G-explicit} after two polynomial integrations.  Cases with one gap zero
follow by continuity.
\end{proof}

For later use, if $\gamma\ge0$, then
\begin{equation}\label{eq:G-minus1gamma}
 \mathfrak G(-1,\gamma)
 =\frac{2(2+3\gamma+2\gamma^2)}{15(1+\gamma)}
 =\frac4{15}+\frac{2\gamma(1+2\gamma)}{15(1+\gamma)}
 \ge\frac4{15}.
\end{equation}

\subsection{Normalized prime-valuation slope vectors}

Recall
\[
 a=v_2(M),\quad b=v_3(A),\quad c=v_2(A),\quad d=v_3(M).
\]
For each prime, define a normalized slope vector $c_p\in\R^2$ by
\begin{align}
 c_2&=\left(-1+\frac a\beta,\ \frac{c}{\beta\vartheta}\right),
 \label{eq:c2}\\
 c_3&=\left(\frac{\vartheta d}{\beta},\ -1+\frac b\beta\right),
 \label{eq:c3}\\
 c_p&=\left(
  \frac{v_p(M)\log p}{\beta\log2},
  \frac{v_p(A)\log p}{\beta\log3}
 \right),\qquad p\ne2,3.
 \label{eq:cp-external}
\end{align}
The logarithmic factorizations of $M$ and $A$ give the conservation law
\begin{equation}\label{eq:slope-conservation}
  \sum_p c_p=(0,0).
\end{equation}
Indeed, the first coordinate sum is
\[
 -1+\frac{a\log2+d\log3+
   \sum_{p\ne2,3}v_p(M)\log p}{\beta\log2}=0,
\]
and similarly for the second coordinate.

The meaning of these vectors can be read directly from the valuation matrices.  Put
$s_k=\beta k/T$.  Up to $o(L_N)$ after division by $T$, the logarithmic $p$-adic cost of the
entry indexed by a scaled column $(x,z)$ is
\begin{center}
\begin{tabular}{@{}ll@{}}
\toprule
Prime & Cost divided by $T L_N\log2$\\
\midrule
$2$ & $x+s_k\,c_2\mathbin{\cdot}(x,z)$\\
$3$ & $z+s_k\,c_3\mathbin{\cdot}(x,z)$\\
$p\ne2,3$ & $s_k\,c_p\mathbin{\cdot}(x,z)$\\
\bottomrule
\end{tabular}
\end{center}
For example,
\[
 v_2(m_k^ua_k^v)=u\bigl(T-\lfloor k\beta\rfloor+ak\bigr)+vck,
\]
and substituting $u=L_Nx$, $v=L_Nz/\vartheta$, $k\sim Ts_k/\beta$ gives the first row.
The $O(u)$ errors from the floors sum to $O(NL_N)=o(H_T)$.

\subsection{The exact continuum limit}

\begin{theorem}[All-prime tropical fraction for the mixed block determinant]
\label{thm:mixed-gini}
Under a hypothetical primitive counterexample, the determinant \eqref{eq:mixed-D} satisfies
\begin{equation}\label{eq:mixed-gini-limit}
 \frac{\log Q_T^{\mathrm{trop}}}{H_T}
 \longrightarrow
 1-\frac38\sum_p\mathfrak G(c_p),
\end{equation}
where the finitely many slope vectors $c_p$ are given by
\eqref{eq:c2}--\eqref{eq:cp-external}.
\end{theorem}

\begin{proof}
The row parameters $s_k=\beta k/T$ become uniformly distributed on $[0,1]$, because
$K_T\sim T/\beta$.  The scaled columns become uniformly distributed on $\DeltaTri$ by
\eqref{eq:semigroup-count}.  The finite rearrangement inequality and Riemann-sum convergence
therefore show that the assignment minimum at prime $p$, after normalization by
$TL_NN\log2$, tends to its constant-column mean plus $\Phi(c_p)$.

Only $p=2$ has constant part $x$, and only $p=3$ has constant part $z$.  Their combined mean is
$2/3$.  By Lemma~\ref{lem:gini-assignment}, the sum of the slope contributions is
\[
 \sum_p\Phi(c_p)
 =\frac16\sum_p(c_{p,1}+c_{p,2})
  -\frac14\sum_p\mathfrak G(c_p).
\]
The first sum vanishes by \eqref{eq:slope-conservation}.  Hence
\[
 \frac{\log Q_T^{\mathrm{trop}}}{TL_NN\log2}
 \longrightarrow
 \frac23-\frac14\sum_p\mathfrak G(c_p).
\]
On the other hand, \eqref{eq:H-asymp} gives
$H_T/(TL_NN\log2)\to2/3$.  Dividing proves
\eqref{eq:mixed-gini-limit}.
\end{proof}

This theorem is exact at the level of minimum Leibniz-term valuations.  It replaces a family
of separate $p$-adic estimates by one finite norm sum satisfying a conservation law.

The four base and cross depths already give a useful branch-sensitive bound without knowing
any external prime factorization.

\begin{corollary}[Depth-sensitive three-vector bound]\label{cor:three-vector}
For every hypothetical primitive counterexample,
\begin{equation}\label{eq:three-vector}
 \limsup_{T\to\infty}
 \frac{\log Q_T^{\mathrm{trop}}}{H_T}
 \le
 1-\frac38\Bigl(
   \mathfrak G(c_2)+\mathfrak G(c_3)
   +\mathfrak G(-c_2-c_3)\Bigr).
\end{equation}
Equality in this upper bound can occur only when all nonzero external slope vectors are
positively collinear.
\end{corollary}

\begin{proof}
The external vectors satisfy
$\sum_{p\ne2,3}c_p=-c_2-c_3$.  Hence the triangle inequality gives
\[
 \sum_{p\ne2,3}\mathfrak G(c_p)
 \ge\mathfrak G(-c_2-c_3).
\]
Substitute this in Theorem~\ref{thm:mixed-gini}.  Equality in a norm triangle inequality
requires all summands to lie on one nonnegative ray.
\end{proof}

The equality condition detects exactly the second-iterate kernel branch reviewed in the source
report.  Indeed the external vectors $c_p$ are rescaled versions of
$(v_p(M),v_p(A))$.  They are all collinear if and only if the two external valuation vectors
of $M$ and $A$ are rationally dependent, equivalently if the kernel $K_\beta$ is nonzero.
Thus:

\begin{corollary}[The tropical norm sees the kernel dichotomy]\label{cor:gini-kernel}
In the multiplicatively independent branch $K_\beta=0$, the inequality in
\eqref{eq:three-vector} is strict.  In the dependent branch $K_\beta\ne0$, the external
triangle inequality is an equality (including the axis cases in which one output is
$\{2,3\}$-smooth).  For fixed $M,A$, the strict gap in the independent branch is explicit:
\begin{equation}\label{eq:external-angular-defect}
 \Delta_{\rm ext}
 =\sum_{p\ne2,3}\mathfrak G(c_p)
  -\mathfrak G\left(\sum_{p\ne2,3}c_p\right)>0.
\end{equation}
\end{corollary}

There is no uniform positive lower bound for $\Delta_{\rm ext}$: distinct rational rays can
approach one another as the valuations grow.  Nevertheless, the quantity is a computable
branch invariant and gives a sharper determinant threshold for every fixed hypothetical
primitive pair.

\subsection{The universal \texorpdfstring{$4/5$}{4/5} ceiling}

\begin{theorem}[Sharp universal min-plus ceiling]\label{thm:four-fifths}
For every hypothetical primitive counterexample,
\begin{equation}\label{eq:G-lower}
  \sum_p\mathfrak G(c_p)\ge\frac8{15}.
\end{equation}
Consequently
\begin{equation}\label{eq:four-fifths}
 \limsup_{T\to\infty}
 \frac{\log Q_T^{\mathrm{trop}}}{H_T}
 \le\frac45.
\end{equation}
The constant $4/5$ is the sharp supremum allowed by the normalized valuation constraints,
although equality would require a degenerate boundary incompatible with a genuine
nonintegral solution.
\end{theorem}

\begin{proof}
By the primitive theorem, either $a=0$ or $b=0$.  Suppose first that $a=0$.  Then
$c_2=(-1,\gamma)$ with $\gamma=c/(\beta\vartheta)\ge0$.  By
\eqref{eq:G-minus1gamma}, $\mathfrak G(c_2)\ge4/15$.  Conservation and the triangle inequality
for the norm give
\[
  \sum_{p\ne2}\mathfrak G(c_p)
  \ge\mathfrak G\!\left(\sum_{p\ne2}c_p\right)
  =\mathfrak G(-c_2)
  =\mathfrak G(c_2).
\]
Thus the total is at least $8/15$.  If $b=0$, apply the same argument to
$c_3=(\eta,-1)$, where $\eta=\vartheta d/\beta\ge0$.  Substitution in
\eqref{eq:mixed-gini-limit} gives
$1-(3/8)(8/15)=4/5$.
\end{proof}

The proof is short because the difficult arithmetic has been compressed into two facts: the
slope vectors sum to zero, and one structural vector reaches from $-1$ to a nonnegative
vertex.  No averaging estimate alone reveals this geometry.

\subsection{Cross-coprime constants: \texorpdfstring{$7/10$ and $3/10$}{7/10 and 3/10}}

Assume the stronger branch condition
\begin{equation}\label{eq:cross-coprime}
  \gcd(MA,6)=1.
\end{equation}
Then $a=b=c=d=0$,
\[
  c_2=(-1,0),
  \qquad
  c_3=(0,-1),
  \qquad
  \sum_{p\ne2,3}c_p=(1,1).
\]

\begin{corollary}[All-prime cross-coprime ceiling]\label{cor:seven-tenths}
Under \eqref{eq:cross-coprime},
\begin{equation}\label{eq:seven-tenths}
 \limsup_{T\to\infty}
 \frac{\log Q_T^{\mathrm{trop}}}{H_T}
 \le\frac7{10}.
\end{equation}
If the external normalized valuation vectors are not all positively collinear, the inequality
is strict by a branch-dependent positive amount.
\end{corollary}

\begin{proof}
By \eqref{eq:G-special}, the structural vectors each contribute $4/15$.  The triangle
inequality gives
\[
  \sum_{p\ne2,3}\mathfrak G(c_p)
  \ge\mathfrak G(1,1)=\frac4{15}.
\]
Thus the total norm is at least $4/5$, and
$1-(3/8)(4/5)=7/10$.  Equality in the triangle inequality requires positive collinearity.
\end{proof}

The places $2$ and $3$ alone have an even smaller exact continuum contribution.

\begin{corollary}[The structural-place $3/10$ constant]\label{cor:three-tenths}
Under \eqref{eq:cross-coprime}, let
$Q_T^{(2,3)}=2^{\tau_{2,T}}3^{\tau_{3,T}}$.  Then
\begin{equation}\label{eq:three-tenths}
  \frac{\log Q_T^{(2,3)}}{H_T}\longrightarrow\frac3{10}.
\end{equation}
\end{corollary}

\begin{proof}
For $c=(-1,0)$,
\[
 \Phi(c)=\frac{-1}{6}-\frac14\cdot\frac4{15}=-\frac7{30},
\]
and the same holds for $(0,-1)$.  Adding the constant mean $2/3$ gives
$2/3-7/15=1/5$.  Dividing by the height mean $2/3$ gives $3/10$.
\end{proof}

Equivalently, one can compute this constant from quantiles.  The descending quantile of
$X_1$ on $\DeltaTri$ is $1-\sqrt t$, so
\[
 \int_0^1t(1-\sqrt t)\,dt=\frac1{10},
 \qquad
 \mathbb EX_1=\frac13.
\]
The ratio is $(1/10)/(1/3)=3/10$ in each structural coordinate.

\subsection{Correction and calibration of a later repository note}

A short $p$-adic priority note added to the repository after the attached unified report makes
the correct strategic observation that minimum Leibniz-term valuations leave a constant
fraction of the archimedean height uncovered.  Two details require correction.

First, the displayed structural valuation formulas in that note simultaneously omit the
terms $kuv_2(M)$ and $kvv_3(A)$; they therefore assume both $a=0$ and $b=0$.  The inherited
primitive theorem gives only $\min(a,b)=0$.  Equations
\eqref{eq:c2}--\eqref{eq:c3} retain all four cross/base depths.

Second, the proposed all-prime $2/3$ ceiling is not a valid universal consequence of
``minimum $\le$ average.''  For an external positive weight proportional to $X_1+X_2$, the
minimum assignment is $4/5$ of its average assignment, not at most $2/3$.  The exact Gini
formula shows that the universal normalized ceiling is $4/5$, while the cross-coprime
all-prime ceiling is $7/10$.  The note's qualitative conclusion survives---entrywise
valuation is insufficient---but the constants and branch hypotheses are materially different.

\begin{statusbox}{\statusnew}
The correction is not based on a numerical counterexample to the Alaoglu--Erd\H{o}s
conjecture.  It is an exact audit of the assignment inequality used in the method.  Rational
external valuation directions can approximate the equality configurations arbitrarily well,
so no stronger uniform constant follows merely from integrality of the valuation vectors.
\end{statusbox}

\subsection{Where extra divisibility can occur}

The tropical minimum equals the true determinant valuation whenever its minimizing
permutation is unique.

\begin{proposition}[Unique tropical term forbids cancellation]\label{prop:unique-tropical}
Let $D=\det(E_{ij})$ be a nonzero integer determinant.  Fix a prime $p$, and suppose a unique
permutation $\sigma_0$ minimizes
$\sum_i v_p(E_{i,\sigma(i)})$.  Then
\[
  v_p(D)=\sum_i v_p(E_{i,\sigma_0(i)}).
\]
\end{proposition}

\begin{proof}
Divide the Leibniz expansion by the corresponding power of $p$.  The term indexed by
$\sigma_0$ is a $p$-adic unit, while every other term is divisible by $p$.  The sum is
therefore a unit modulo $p$.
\end{proof}

Consequently, the missing $20\%$ or more can arise only from primes at which the assignment
problem has several minimizers.  The tie geometry is explicit.

\begin{itemize}
\item For $p\ne2,3$, columns tie when
      $v_p(M)(u-u')+v_p(A)(v-v')=0$, so ties lie on rational lines orthogonal to the external
      valuation vector.
\item In the cross-coprime branch, the $2$-adic cost depends only on $u$ and the $3$-adic cost
      only on $v$.  Every vertical or horizontal column fiber creates many equal-cost
      permutations.  Any leading extra valuation must be produced inside these fibers.
\item Cross valuations $v_2(A)$ and $v_3(M)$ tilt the fibers.  The floor term in the row
      profile perturbs the exact affine ordering but does not change the continuum slope
      vector.
\end{itemize}

This turns ``look for $p$-adic cancellation'' into a concrete algebra problem: factor or
control the residual determinants inside equal-slope column fibers, and show that their
valuation supplies at least the missing Gini deficit.


\section{The complete triangular grid and a \texorpdfstring{$\pi/4$}{pi/4} transport ceiling}
\label{sec:fullgrid}

The one-block determinant deliberately uses the exact $O(T)$ population in a unit interval of
exponent space.  The source report also suggests exploiting the full two-dimensional
coordinate grid rather than sorting only by magnitude.  Conditional on a counterexample, that
proposal gives $\asymp T^2$ genuine integer rows below exponent height $T$, exactly the number
that the mixed-semigroup interpolation argument seemed to be missing.  It is therefore
important to determine whether entrywise $p$-adic divisibility becomes strong enough on this
larger node set.

The answer is again negative, but for a different and pleasantly universal reason.  Once all
places are summed, the prime-by-prime assignment minima are bounded by one antimonotone
transport problem for the logarithmic entry matrix.  Both the row heights and the balanced
column weights have the same triangular distribution, and the ratio between antimonotone and
monotone transport is exactly $\pi/4$.

\subsection{Rows and columns}

For $T>0$ define the complete conditional row set
\begin{equation}\label{eq:grid-rows}
 \mathcal R_T=
 \{(n,k)\in\Nzero^2:n+k\beta\le T\},
 \qquad N_T=\#\mathcal R_T.
\end{equation}
For $i=(n_i,k_i)\in\mathcal R_T$ put
\begin{equation}\label{eq:grid-row-values}
 x_i=n_i+k_i\beta,
 \qquad m_i=2^{n_i}M^{k_i}=2^{x_i},
 \qquad a_i=3^{n_i}A^{k_i}=3^{x_i}.
\end{equation}
Because $\beta$ is irrational, the $x_i$ are pairwise distinct.

Let $\mathcal C_T$ be the first $N_T$ points $(u,v)\in\Nzero^2$ in increasing order of
\[
 \lambda=u+\vartheta v,
 \qquad
 y=\log(2^u3^v)=\lambda\log2.
\]
Write $L_T$ for the largest selected $\lambda$.  Standard lattice-point estimates for the two
triangles give
\begin{equation}\label{eq:grid-cardinality}
 N_T=\frac{T^2}{2\beta}+O_\beta(T),
 \qquad
 L_T=\sqrt{2\vartheta N_T}+O_\vartheta(1)
      =T\sqrt{\frac{\vartheta}{\beta}}+O_\beta(1).
\end{equation}

The associated square determinant is
\begin{equation}\label{eq:grid-determinant}
 D_T^{\triangle}
 =\det\!\left(m_i^{u_j}a_i^{v_j}\right)_{i,j=1}^{N_T}
 =\det\!\left(e^{x_i y_j}\right)_{i,j=1}^{N_T}.
\end{equation}
Every entry is an integer.  Ordering the $x_i$ and $y_j$ increasingly, strict total positivity
of the exponential kernel makes $D_T^{\triangle}>0$; equivalently, this is the same
exponential-polynomial zero theorem used by the generalized-Vandermonde part of the source
report.

\begin{statusbox}{\statusnew{} for the asymptotic transport statements below.  Integrality and
nonvanishing are inherited from the mixed-semigroup determinant architecture
\statuslean.}
\end{statusbox}

\subsection{The common triangular distribution}

The normalized row heights $h_i=x_i/T$ and column weights $\ell_j=y_j/(L_T\log2)$ both live in
$[0,1]$.  Their empirical measures have the same limit.

\begin{lemma}[Triangular height law]\label{lem:triangular-law}
For every continuous $f:[0,1]\to\R$,
\begin{align}
 \frac1{N_T}\sum_{i=1}^{N_T}f(h_i)
   &\longrightarrow \int_0^1 2t f(t)\,dt, \label{eq:row-law}\\
 \frac1{N_T}\sum_{j=1}^{N_T}f(\ell_j)
   &\longrightarrow \int_0^1 2t f(t)\,dt. \label{eq:column-law}
\end{align}
Thus the common limiting distribution has cumulative distribution function $F(t)=t^2$ and
quantile function
\begin{equation}\label{eq:triangular-quantile}
 Q(s)=\sqrt s,\qquad 0\le s\le1.
\end{equation}
\end{lemma}

\begin{proof}
After scaling $(n,k)$ to $(n/T,\beta k/T)$, the row lattice becomes equidistributed in the
unit right triangle
\[
 \{(r,s):r,s\ge0,\ r+s\le1\}.
\]
The subtriangle on which $r+s\le t$ has area $t^2/2$, while the whole triangle has area
$1/2$.  Hence the height $r+s$ has distribution function $t^2$.  Boundary discrepancies are
$O(T)$ against $N_T\asymp T^2$, proving \eqref{eq:row-law}.

For the columns, scale $(u,v)$ to $(u/L_T,\vartheta v/L_T)$.  The same unit triangle appears,
and $\ell=(u+\vartheta v)/L_T$ is again the sum of the scaled coordinates.  The boundary shell
created by taking the first exactly $N_T$ points has $O(L_T)$ points against
$N_T\asymp L_T^2$, which proves \eqref{eq:column-law}.
\end{proof}

\subsection{Summing all primewise assignment minima}

For each prime $p$, let
\begin{equation}\label{eq:grid-tau}
 \tau_{p,T}
 =\min_{\sigma\in S_{N_T}}
   \sum_{i=1}^{N_T}
   v_p\!\left(m_i^{u_{\sigma(i)}}a_i^{v_{\sigma(i)}}\right),
 \qquad
 Q_T^{\triangle}=\prod_p p^{\tau_{p,T}}.
\end{equation}
As before, $Q_T^{\triangle}$ divides every Leibniz term and hence divides
$D_T^{\triangle}$.  The key observation is that taking a separate minimizing permutation at
each prime can only decrease the sum relative to using one permutation at all primes.

\begin{lemma}[All-place transport domination]\label{lem:all-place-transport}
For every $T$,
\begin{equation}\label{eq:all-place-min}
 \log Q_T^{\triangle}
 \le
 \min_{\sigma\in S_{N_T}}
 \sum_{i=1}^{N_T}x_i y_{\sigma(i)}.
\end{equation}
\end{lemma}

\begin{proof}
For a permutation $\sigma$, put
\[
 C_p(\sigma)=
 \sum_i v_p\!\left(m_i^{u_{\sigma(i)}}a_i^{v_{\sigma(i)}}\right)\log p.
\]
Then
\[
 \sum_p\min_\sigma C_p(\sigma)
 \le \min_\sigma\sum_p C_p(\sigma).
\]
Unique factorization and \eqref{eq:grid-row-values} give, entry by entry,
\[
 \sum_p v_p(m_i^{u_j}a_i^{v_j})\log p
 =\log(m_i^{u_j}a_i^{v_j})
 =x_i y_j.
\]
Substitution proves \eqref{eq:all-place-min}.
\end{proof}

Notice how little arithmetic remains in the last inequality: it is valid in every valuation
branch, includes all primes, and does not use the primitive-depth normalization.  The price is
that it is only an upper bound for the common termwise divisor.  That is exactly what is needed
for a no-go theorem.

\subsection{Optimal transport and the \texorpdfstring{$\pi/4$}{pi/4} ratio}

Let
\begin{align}
 E_T^-&=\min_{\sigma\in S_{N_T}}\sum_i x_i y_{\sigma(i)},
 \label{eq:grid-Eminus}\\
 E_T^+&=\max_{\sigma\in S_{N_T}}\sum_i x_i y_{\sigma(i)}.
 \label{eq:grid-Eplus}
\end{align}
The rearrangement inequality pairs the two ordered lists oppositely for $E_T^-$ and in the
same order for $E_T^+$.

\begin{theorem}[Triangular-grid transport constants]\label{thm:grid-transport}
As $T\to\infty$,
\begin{align}
 \frac{E_T^-}{N_TTL_T\log2}
   &\longrightarrow
     \int_0^1\sqrt{t(1-t)}\,dt
      =\frac\pi8, \label{eq:grid-min-limit}\\
 \frac{E_T^+}{N_TTL_T\log2}
   &\longrightarrow
     \int_0^1t\,dt
      =\frac12. \label{eq:grid-max-limit}
\end{align}
Consequently
\begin{equation}\label{eq:grid-ratio}
 \frac{E_T^-}{E_T^+}\longrightarrow\frac\pi4
 =0.7853981633\ldots .
\end{equation}
\end{theorem}

\begin{proof}
By Lemma~\ref{lem:triangular-law}, the increasing empirical quantile functions of the two
normalized lists converge in $L^1$ to $Q(t)=\sqrt t$.  The discrete rearrangement sums are
Riemann sums for the monotone and antimonotone couplings.  Thus
\[
 \frac{E_T^-}{N_TTL_T\log2}
 \longrightarrow\int_0^1Q(t)Q(1-t)\,dt,
 \qquad
 \frac{E_T^+}{N_TTL_T\log2}
 \longrightarrow\int_0^1Q(t)^2\,dt.
\]
The first integral is the beta integral
$B(3/2,3/2)=\Gamma(3/2)^2/\Gamma(3)=\pi/8$; the second is $1/2$.
\end{proof}

The determinant has $N_T!$ Leibniz terms, each at most $e^{E_T^+}$.
Since $N_T\asymp T^2$ while $TL_TN_T\asymp T^4$,
\begin{equation}\label{eq:grid-factorial-lower-order}
 \log(N_T!)=O(T^2\log T)=o(N_TTL_T).
\end{equation}
Combining the preceding statements gives the promised ceiling.

\begin{corollary}[Full-grid termwise-divisor ceiling]\label{cor:grid-pi-four}
With the notation above,
\begin{equation}\label{eq:grid-pi-four}
 \limsup_{T\to\infty}
 \frac{\log Q_T^{\triangle}}
      {E_T^++\log(N_T!)}
 \le\frac\pi4.
\end{equation}
Thus the complete $\asymp T^2$ row grid does not make the all-prime common divisor exceed the
sharp largest-Leibniz-term archimedean bound.  Entrywise valuation alone leaves at least the
fraction $1-\pi/4=0.2146018\ldots$ uncovered at leading order.
\end{corollary}

\begin{proof}
Lemma~\ref{lem:all-place-transport} gives
$\log Q_T^{\triangle}\le E_T^-$.  Divide by the standard determinant upper bound
$E_T^++\log(N_T!)$ and apply Theorem~\ref{thm:grid-transport} and
\eqref{eq:grid-factorial-lower-order}.
\end{proof}

\begin{warning}[What the $\pi/4$ theorem does not exclude]\label{warn:grid-actual}
The denominator in \eqref{eq:grid-pi-four} is the largest-term upper bound, not the true size
of $D_T^{\triangle}$.  Generalized Vandermonde determinants exhibit substantial alternating
cancellation.  A sufficiently sharp archimedean estimate could lower the relevant budget
below $E_T^-$, in which case the termwise divisor might become decisive.  The theorem excludes
the naive ``more rows plus the same Hadamard/Leibniz estimate'' route; it does not exclude a
joint $p$-adic/free-energy argument.
\end{warning}

\subsection{The Harish--Chandra--Itzykson--Zuber identity}

The exact analytic object behind that warning is classical.  For diagonal Hermitian matrices
$X=\operatorname{diag}(x_1,\ldots,x_N)$ and
$Y=\operatorname{diag}(y_1,\ldots,y_N)$ with distinct increasing entries, normalized Haar
measure on the unitary group satisfies
\begin{equation}\label{eq:hciz}
 \int_{U(N)}
   \exp\!\bigl(\operatorname{Tr}(XUYU^*)\bigr)\,dU
 =\frac{\prod_{r=1}^{N-1}r!}{\Delta(x)\Delta(y)}
   \det(e^{x_i y_j})_{i,j=1}^{N},
\end{equation}
where
$\Delta(x)=\prod_{i<j}(x_j-x_i)$ and similarly for $\Delta(y)$
\cite{HarishChandra1957,ItzyksonZuber1980,McSwiggen2018}.  Hence
\begin{equation}\label{eq:hciz-det}
 \log D_T^{\triangle}
 =\log\Delta(x)+\log\Delta(y)
  -\sum_{r=1}^{N_T-1}\log(r!)
  +\log\mathcal I_T,
\end{equation}
with $\mathcal I_T$ the unitary integral in \eqref{eq:hciz}.

The large powers of $T$ hidden separately in the two Vandermonde factors and the factorial
product cancel at leading logarithmic scale.  What remains is a genuine free-energy problem
for two triangular empirical measures.  The Baker separation theorem of
Section~\ref{sec:baker} gives a polynomial lower bound for the spacings in $\Delta(x)$, so the
row Vandermonde cannot collapse through anomalously close conditional exponents.  What is not
known is a bound on the full combination \eqref{eq:hciz-det} strong enough to beat the
primewise divisor.

\begin{target}[HCIZ--valuation comparison]\label{target:hciz}
Determine the $T^4$-scale limit, or a sufficiently sharp upper bound, of
\[
 \log D_T^{\triangle}-\log Q_T^{\triangle}
\]
for the conditional triangular row lattice and balanced mixed columns.  It would suffice to
prove that this quantity is negative for all sufficiently large $T$.  Corollary
\ref{cor:grid-pi-four} shows that such a result must use either the HCIZ/Vandermonde
cancellation in \eqref{eq:hciz-det}, extra $p$-adic cancellation beyond the assignment minima,
or both.
\end{target}

This formulation has a useful advantage over a generic interpolation determinant: both
limiting measures and all lattice densities are explicit.  It is therefore suitable for
numerical free-energy reconnaissance before attempting a proof.

\section{The first tied \texorpdfstring{$p$}{p}-adic layer is lower order}
\label{sec:firstlayer}

Theorem~\ref{thm:four-fifths} says that extra divisibility must come from ties among
minimum-valuation Leibniz terms.  Ties are especially abundant in the cross-coprime branch
\eqref{eq:cross-coprime}: at $p=2$ all columns with the same $u$ have the same valuation
profile, and at $p=3$ all columns with the same $v$ do.  It is natural to hope that cancellation
inside those large fibers supplies the missing leading-order fraction.

The present section computes that first cancellation layer exactly.  It factors into ordinary
Vandermondes of the odd row parts.  Standard upper bounds for $p$-adic linear forms in two
logarithms then show that the whole first layer is only
$O(T^{3/2}(\log T)^C)$, whereas the mixed determinant has height $\asymp T^{5/2}$.  This does
not bound the final determinant valuation---higher tropical layers may cancel again---but it
rules out a one-step fiber explanation of the missing fraction.

\subsection{Tropical initial form at two}

Assume throughout this section that $\gcd(MA,6)=1$.  For the block rows, write
\begin{equation}\label{eq:odd-output}
 q_k=a_{T,k}=3^{n_k}A^k,
 \qquad n_k=T-\lfloor k\beta\rfloor.
\end{equation}
Both $M$ and every $q_k$ are odd, and
\begin{equation}\label{eq:two-entry-factor}
 m_{T,k}^{u}a_{T,k}^{v}
 =2^{u n_k}\,M^{ku}q_k^v.
\end{equation}
For each integer $u\ge0$, let
\[
 V_u=\{v:(u,v)\text{ is one of the selected }N\text{ columns}\},
 \qquad s_u=\#V_u.
\]
Because columns are selected by the inequality $u+\vartheta v\le L_N$ up to one boundary
shell, each nonempty $V_u$ is an initial interval of integers.  Sort the distinct $u$ values
in increasing order.  Since $n_0>n_1>\cdots>n_{N-1}$, every minimizing assignment for the
cost $un_k$ sends one consecutive row block $R_u$ of size $s_u$ to the columns in the fiber
$\{u\}\times V_u$.

Let $\tau_{2,T}$ be the minimum from \eqref{eq:taup}.  The sum of all Leibniz terms having
exactly this valuation is the $2$-adic tropical initial form.  Up to an overall sign and an odd
factor, it is
\begin{equation}\label{eq:initial-two}
 I_{2,T}
 =\prod_u
   \det\!\left(q_k^v\right)_{k\in R_u,\ v\in V_u}
 =\prod_u\prod_{\substack{i<j\\i,j\in R_u}}(q_j-q_i).
\end{equation}
The first equality follows by grouping the optimal permutations by fibers and factoring the
odd number $M^{ku}$ from each row in a fixed $u$ block; the second is the ordinary Vandermonde
identity because $V_u$ is consecutive.

\begin{proposition}[Exact first $2$-adic initial form]\label{prop:initial-two}
There is an integer $Z_{2,T}$ and an integer gap $g_{2,T}\ge\lfloor\beta\rfloor$ such that
\begin{equation}\label{eq:two-layer-expansion}
 2^{-\tau_{2,T}}D_T
 =\varepsilon_{2,T} I_{2,T}+2^{g_{2,T}}Z_{2,T},
\end{equation}
where $\varepsilon_{2,T}$ is odd.  Consequently
\begin{equation}\label{eq:two-layer-min}
 v_2(D_T)-\tau_{2,T}
 \ge \min\{v_2(I_{2,T}),g_{2,T}\},
\end{equation}
and equality holds whenever the two numbers inside the minimum are unequal.
\end{proposition}

\begin{proof}
The rearrangement inequality with repeated column weights identifies the minimizing
permutations exactly as above, and summing their signed unit parts gives
\eqref{eq:initial-two}.  Any nonminimizing assignment must move at least one column across two
adjacent $u$ fibers.  An exchange across the corresponding boundary changes the cost by
\[
 (u'-u)(n_i-n_j).
\]
Here $u'-u\ge1$, while consecutive row depths differ by
$\lfloor(k+1)\beta\rfloor-\lfloor k\beta\rfloor\ge\lfloor\beta\rfloor$.  Hence every
nonminimal Leibniz term has valuation at least $\tau_{2,T}+\lfloor\beta\rfloor$, proving the
expansion.  The valuation statement follows from the ultrametric inequality.
\end{proof}

\subsection{The symmetric initial form at three}

Put
\begin{equation}\label{eq:three-unit-input}
 r_k=m_{T,k}=2^{n_k}M^k.
\end{equation}
At $p=3$ the entry factors as
\[
 m_{T,k}^ua_{T,k}^v
 =3^{vn_k}\,r_k^uA^{kv}.
\]
For a selected value of $v$, let $U_v$ be the consecutive set of selected $u$ values and let
$R'_v$ be the corresponding consecutive block of rows in the minimum assignment.  The same
argument gives
\begin{equation}\label{eq:initial-three}
 I_{3,T}
 =\prod_v\prod_{\substack{i<j\\i,j\in R'_v}}(r_j-r_i)
\end{equation}
and
\begin{equation}\label{eq:three-layer-expansion}
 3^{-\tau_{3,T}}D_T
 =\varepsilon_{3,T}I_{3,T}+3^{g_{3,T}}Z_{3,T},
 \qquad g_{3,T}\ge\lfloor\beta\rfloor,
\end{equation}
with $\varepsilon_{3,T}$ a $3$-adic unit.

\subsection{Two-logarithm bounds inside the fibers}

For $i<j$, set $d=j-i$ and
$h=n_i-n_j=\lfloor j\beta\rfloor-\lfloor i\beta\rfloor$.  The common factors in the row
values are units at the relevant place, so
\begin{align}
 v_2(q_j-q_i)&=v_2(A^d-3^h), \label{eq:qdiff-two}\\
 v_3(r_j-r_i)&=v_3(M^d-2^h). \label{eq:rdiff-three}
\end{align}
Neither difference vanishes: $A^d=3^h$ or $M^d=2^h$ would give $d\beta=h$, contrary to the
irrationality of $\beta$.

A standard consequence of $p$-adic linear-form theory is that for fixed multiplicatively
independent positive rationals $\alpha_1,\alpha_2$ and a fixed prime $p$, there are effectively
computable constants $C_0,C_1$ such that
\begin{equation}\label{eq:padic-two-log}
 v_p(\alpha_1^d-\alpha_2^h)
 \le C_0+C_1\bigl(\log(2+\max\{d,h\})\bigr)^2.
\end{equation}
Much sharper versions, with linear dependence on $\log B$, are known in important two-logarithm
regimes; the square is more than sufficient here
\cite{BugeaudLaurent1996,Bugeaud2002,Yu1998}.  Apply this to $(A,3,p=2)$ and
$(M,2,p=3)$.  Since $d,h=O_\beta(N)$, each factor in
\eqref{eq:initial-two}--\eqref{eq:initial-three} has valuation
$O_{M,A}((\log N)^2)$.

The number of pairs within the vertical fibers is
\begin{equation}\label{eq:vertical-pairs}
 \sum_u\binom{s_u}{2}
 =\frac{L_N^3}{6\vartheta^2}+O_\vartheta(L_N^2)
 =O_\vartheta(N^{3/2}),
\end{equation}
while the number within the horizontal fibers is
\begin{equation}\label{eq:horizontal-pairs}
 \sum_v\binom{\#U_v}{2}
 =\frac{L_N^3}{6\vartheta}+O_\vartheta(L_N^2)
 =O_\vartheta(N^{3/2}).
\end{equation}
We obtain the following rigorous scale separation.

\begin{theorem}[First tied layer is lower order]\label{thm:first-layer-small}
In the cross-coprime branch,
\begin{equation}\label{eq:first-layer-small}
 v_2(I_{2,T})+v_3(I_{3,T})
 =O_{M,A}\!\left(N^{3/2}(\log N)^2\right)
 =O_{M,A}\!\left(T^{3/2}(\log T)^2\right).
\end{equation}
In particular,
\begin{equation}\label{eq:first-layer-negligible}
 \frac{v_2(I_{2,T})\log2+v_3(I_{3,T})\log3}{H_T}
 \longrightarrow0.
\end{equation}
\end{theorem}

\begin{proof}
Sum \eqref{eq:padic-two-log} over the pairs in
\eqref{eq:vertical-pairs} and \eqref{eq:horizontal-pairs}, using
\eqref{eq:qdiff-two}--\eqref{eq:rdiff-three}.  Finally use
$N\asymp_\beta T$ and $H_T\asymp_\beta T^{5/2}$.
\end{proof}

\begin{statusbox}{\statusnew{} modulo the quoted classical $p$-adic linear-form estimate
\eqref{eq:padic-two-log}.  The factorization of the tropical initial forms is elementary and
is a particularly attractive Lean target.}
\end{statusbox}

\subsection{What a leading excess would have to do}

Theorem~\ref{thm:first-layer-small} does not prove
$v_2(D_T)-\tau_{2,T}=o(H_T)$ or its $3$-adic analogue.  If the initial form is divisible by a
power exceeding the first tropical gap, terms from the next valuation layer enter; those can
cancel, exposing another layer, and so on.  A leading-order excess is therefore possible only
through a \emph{cascade} of cancellations across a number of tropical layers growing with
$T$.  Cancellation confined to the minimum-permutation fibers is quantitatively negligible.

This yields a more precise version of the $p$-adic research target.

\begin{target}[No shallow-layer proof]\label{target:cascade}
Either prove that the tropical cancellation cascade has total depth $o(H_T)$---which would
retire the single-block $p$-adic strategy---or identify a structural identity forcing
$\gg H_T$ cumulative cancellation across successively nonminimal assignments.  Pairwise
congruences such as \eqref{eq:qdiff-two} and \eqref{eq:rdiff-three} cannot by themselves
supply the latter, because $p$-adic logarithm bounds make each congruence only polylogarithmic
in the row separation.
\end{target}

\section{Synthesis of the determinant audits}
\label{sec:synthesis}

The preceding calculations put several superficially different determinant proposals into the
same currency: the fraction of the leading archimedean budget supplied automatically by
valuations of the entries.  The numerical constants should not be compared as though the
underlying determinants were identical---the ordinary, mixed-block, and complete-grid systems
have different sizes and different upper bounds---but they do show which mechanism is closest
to a contradiction before any genuine cancellation is used.

\begin{table}[htbp]
\centering
\caption{Leading termwise-divisor ceilings obtained in this continuation.}
\label{tab:ceilings}
\begin{tabularx}{\textwidth}{@{}>{\raggedright\arraybackslash}p{0.29\textwidth}
 >{\raggedright\arraybackslash}p{0.27\textwidth}
 >{\raggedright\arraybackslash}X@{}}
\toprule
System & Ceiling or exact limit & Meaning \\
\midrule
Ordinary input/output Vandermonde product &
$1-\dfrac{\delta_2\log2+\delta_3\log3}{3\beta\log6}
 <0.8710491$ &
Exact fixed-support divisor fraction; all primes already occurring in $M$ and $A$ included.\\[0.7em]
Mixed single block, arbitrary branch &
$1-\dfrac38\sum_p\mathfrak G(c_p)\le4/5$ &
Exact all-prime tropical limit and universal primitive-depth ceiling.\\[0.7em]
Mixed single block, cross-coprime &
$\le7/10$ &
Uses $c_2=(-1,0)$, $c_3=(0,-1)$ and the external sum $(1,1)$.\\[0.7em]
Mixed single block, only structural primes in the cross-coprime branch &
$3/10$ &
Exact contribution of $p=2,3$ before external primes or cancellation.\\[0.7em]
Complete triangular row grid &
$\le\pi/4=0.7853982\ldots$ &
All-prime common divisor versus the sharp largest-Leibniz-term transport upper bound.\\[0.7em]
First tied fibers at $2$ and $3$ &
$o(1)$ of the mixed-block budget &
Their valuations are $O(T^{3/2}(\log T)^2)$ against $H_T\asymp T^{5/2}$.\\
\bottomrule
\end{tabularx}
\end{table}

Three conclusions survive every normalization.

\begin{enumerate}
\item Merely adding every prime in the fixed support does not close any determinant.  It does,
      however, improve the diagnosis substantially: the universal mixed-block ceiling is
      $4/5$, not a heuristic ``about one half''.
\item The exact primitive theorem is quantitatively useful.  The condition
      $\min(v_2(M),v_3(A))=0$ creates one structural slope vector of norm at least $4/15$; the
      conservation law forces a second norm contribution of at least the same size.
\item A successful $p$-adic proof must be a cancellation theorem, not a divisibility-by-every-
      term theorem.  In the easiest branch, even the first large tied family contributes only
      a lower-order amount.  The required phenomenon would have to repeat through many
      tropical layers or interact with a much sharper archimedean free energy.
\end{enumerate}

\subsection{Correction of the post-attachment ``one half / two thirds'' note}
\label{subsec:report7-correction}

After the attached unified report was produced, the repository acquired a short exploratory
note, \lean{Analysis/ExponentialIdentities/Article/new-reports/report-7.md}, at commit
\sha{ff9a9f5aa3a3f02599de93fc053e60918f269855}.  That note correctly identified the distinction
between termwise min-plus divisibility and cancellation beyond the tropical minimum, and its
real/$p$-adic compatibility diagnosis remains valuable.  It proposed, however, a structural
$1/2$ ceiling and an all-prime $2/3$ ceiling.  Neither constant is valid in the stated
generality.

The structural averaging argument omitted the base and cross depths that appear in
\eqref{eq:c2}--\eqref{eq:c3}.  Even under $v_2(M)=0$, the terms
$v_2(A)kv$ and $v_3(M)ku+v_3(A)kv$ do not disappear.  The exact structural contribution is
obtained by retaining $c_2$ and $c_3$ in Theorem~\ref{thm:mixed-gini}; it can exceed one half
of $H_T$.

The all-prime step used the assertion that an oppositely sorted nonnegative triangular weight
has assignment minimum at most two thirds of its mean.  For the weight $Y=X+Z$ on the standard
triangle, the exact values from Lemma~\ref{prop:G-explicit} are
\[
 \mathbb E Y=\frac23,
 \qquad
 \Phi(1,1)=\frac4{15},
 \qquad
 \frac{\Phi(1,1)}{\mathbb EY/2}=\frac45.
\]
Thus $4/5$, not $2/3$, is already attained by the relaxed external rank-one grading.  The
Gini-norm identity explains the general case and leads to the rigorous ceiling in
Theorem~\ref{thm:four-fifths}.  This is a correction of a calculation in an exploratory note,
not a criticism of the source report itself; the attached unified report did not contain the
claim.

\begin{statusbox}{\statusnew{} --- the correction follows from the exact transport formula,
not from numerical experimentation.}
\end{statusbox}


\section{Directed-rounding extension through \texorpdfstring{$2^{27}$}{2\^{}27}}
\label{sec:finite27}

The attached report supplied a complete MPFR scan through $2^{26}$.  Because a finite
extension is independent of the unresolved transcendence barrier and gives a useful stronger
floor for every hypothetical generator, I extracted the literal C listing from its
Appendix~``Directed-rounding verifier'', compiled it against MPFR~4.2.2, replayed an old audit
chunk, and scanned the next full dyadic range.

\subsection{Manifest audit}

The literal extracted source has SHA-256 digest
\begin{center}
 \sha{f888c0d89be9366f01e2d26f339dcbe446556d39d4c5dbcf8f579ec0d0052e1d}.
\end{center}
This does \emph{not} equal the digest
\sha{ec709623f812478fbb7fb5567aa59fbfe3df0ae639ca366bc7c72e278bcfee2f}
printed in the attached report.  None of the ordinary line-ending, final-newline, or tab
normalizations explains the discrepancy.  The report therefore contains a stale or
externally generated source-identity hash.

The semantic reproducibility check succeeds.  Replaying the old interval
$[29{,}360{,}128,33{,}554{,}432)$ at $192$ bits gives exactly the report's records
\[
 \begin{split}
 &\text{lower: }1.4439124872349507\times10^{-19}
   \text{ at }(30{,}103{,}832,713{,}403{,}245{,}915),\\
 &\text{upper: }5.4132120119800561\times10^{-21}
   \text{ at }(29{,}411{,}704,687{,}581{,}893{,}443),
 \end{split}
\]
with zero failures.  Thus the listing used here reproduces the old executable's mathematical
behavior on an audited four-million-input overlap even though the printed file digest does not
identify it.

\begin{statusbox}{\statuscomp{} --- this section is a directed-rounding software result.  It
is stronger numerically than the Lean theorem but is not a kernel proof.  The manifest hash
mismatch is explicitly part of the trust boundary.}
\end{statusbox}

\subsection{The new range}

The half-open interval $[2^{26},2^{27})$ was divided into sixteen chunks of width $2^{22}$.
The executable used $192$-bit MPFR endpoints and outward rounding exactly as in the source
report.  There are $2^{26}$ integers in the range; the first, $2^{26}$, is the sole power of
two and is intentionally skipped.  Hence the required nonpower count is
$2^{26}-1=67{,}108{,}863$.

The concatenated canonical extension log has SHA-256 digest
\begin{center}
 \sha{b11843448cd8946442aba13758b5b666e4ac6eae50d2243fbdc32d83274cea8d}.
\end{center}
Every chunk reports zero failures, and the checked counts sum to the required total.

{\small
\begin{longtable}{@{}r p{0.27\textwidth} r r r@{}}
\caption{Directed-rounding extension on $[2^{26},2^{27})$.}
\label{tab:scan27}\\
\toprule
Chunk & Integer range & Checked & Smallest lower gap & Smallest upper gap \\
\midrule
\endfirsthead
\toprule
Chunk & Integer range & Checked & Smallest lower gap & Smallest upper gap \\
\midrule
\endhead
0 & 67{,}108{,}864--71{,}303{,}167 & 4{,}194{,}303 & $4.2837\times10^{-20}$ & $1.7723\times10^{-19}$ \\
1 & 71{,}303{,}168--75{,}497{,}471 & 4{,}194{,}304 & $2.2961\times10^{-20}$ & $1.1433\times10^{-20}$ \\
2 & 75{,}497{,}472--79{,}691{,}775 & 4{,}194{,}304 & $1.2813\times10^{-19}$ & $1.6063\times10^{-20}$ \\
3 & 79{,}691{,}776--83{,}886{,}079 & 4{,}194{,}304 & $2.5420\times10^{-20}$ & $5.2012\times10^{-20}$ \\
4 & 83{,}886{,}080--88{,}080{,}383 & 4{,}194{,}304 & $2.6166\times10^{-20}$ & $3.8496\times10^{-20}$ \\
5 & 88{,}080{,}384--92{,}274{,}687 & 4{,}194{,}304 & $1.0923\times10^{-19}$ & $5.7144\times10^{-20}$ \\
6 & 92{,}274{,}688--96{,}468{,}991 & 4{,}194{,}304 & $1.1017\times10^{-20}$ & $2.3227\times10^{-20}$ \\
7 & 96{,}468{,}992--100{,}663{,}295 & 4{,}194{,}304 & $3.0759\times10^{-20}$ & $7.4810\times10^{-20}$ \\
8 & 100{,}663{,}296--104{,}857{,}599 & 4{,}194{,}304 & $5.6465\times10^{-20}$ & $6.4744\times10^{-20}$ \\
9 & 104{,}857{,}600--109{,}051{,}903 & 4{,}194{,}304 & $8.8144\times10^{-20}$ & $8.3662\times10^{-21}$ \\
10 & 109{,}051{,}904--113{,}246{,}207 & 4{,}194{,}304 & $3.8892\times10^{-20}$ & $7.0327\times10^{-21}$ \\
11 & 113{,}246{,}208--117{,}440{,}511 & 4{,}194{,}304 & $9.9365\times10^{-21}$ & $2.1814\times10^{-20}$ \\
12 & 117{,}440{,}512--121{,}634{,}815 & 4{,}194{,}304 & $4.4964\times10^{-20}$ & $5.4132\times10^{-21}$ \\
13 & 121{,}634{,}816--125{,}829{,}119 & 4{,}194{,}304 & $1.0774\times10^{-20}$ & $3.4093\times10^{-20}$ \\
14 & 125{,}829{,}120--130{,}023{,}423 & 4{,}194{,}304 & $4.3609\times10^{-20}$ & $8.3073\times10^{-22}$ \\
15 & 130{,}023{,}424--134{,}217{,}727 & 4{,}194{,}304 & $5.9082\times10^{-21}$ & $5.1326\times10^{-20}$ \\
\midrule
\multicolumn{2}{@{}r}{\textbf{total}} & \textbf{67{,}108{,}863} & & \\
\bottomrule
\end{longtable}
}

\subsection{New extremal records and high-precision replay}

The smallest lower logarithmic margin in the extension is
\begin{equation}\label{eq:new-lower-record}
 5.9082259992879571\times10^{-21}
 \quad\text{at}\quad
 (m,n)=(132{,}833{,}765,7{,}501{,}348{,}219{,}569).
\end{equation}
A $120$-decimal-digit reevaluation gives
\[
 m^{\vartheta}-n
 =6.3939754533470409837680125510\ldots\times10^{-8}.
\]

The smallest upper logarithmic margin is the new global record
\begin{equation}\label{eq:new-upper-record}
 8.3073231820333117\times10^{-22}
 \quad\text{at}\quad
 (m,n)=(129{,}799{,}345,7{,}231{,}571{,}292{,}851),
\end{equation}
where
\[
 n+1-m^{\vartheta}
 =8.6669904355818956853297952288\ldots\times10^{-9}.
\]
Both one-input cases were rerun independently at $320$ and $512$ bits.  The directed margins
and nearest integer were unchanged in every displayed digit, and both replays reported zero
failures.

\begin{theorem}[Software-level finite exclusion through $2^{27}$]
\label{thm:finite27}
At the stated MPFR software-trust level, for every integer $1\le m<2^{27}$,
\[
 m^{\vartheta}\in\N
 \quad\Longrightarrow\quad
 m\text{ is a power of two}.
\]
Consequently every nonintegral real $x$ satisfying $2^x,3^x\in\Z$ obeys
\begin{equation}\label{eq:x-gt-27}
 x>27.
\end{equation}
\end{theorem}

\begin{proof}
The attached report's complete scan covers $1\le m<2^{26}$, and the sixteen new chunks cover
$2^{26}\le m<2^{27}$.  In each case the directed logarithmic interval for
$\vartheta\log m$ lies strictly between the intervals for the logarithms of two consecutive
integers, except when $m$ is a power of two.  If $m=2^x$ came from a nonintegral solution and
$m<2^{27}$, then $m^\vartheta=3^x$ would be an integer at a non-power input, contradicting the
scan.  The endpoint $m=2^{27}$ corresponds to the integral exponent $27$, so a nonintegral
solution must have $m>2^{27}$ and hence $x>27$.
\end{proof}

This raises the software-level spacing ratio in the optimized finite-difference comparison
from $26\log3$ to $27\log3=29.662531\ldots$.  It does not change any kernel-verified theorem;
the formal floor remains $x\ge12$ until the finite certificates are replayed in Lean.

\section{Additional proof routes pushed to their failure points}
\label{sec:attempts}

The determinant calculations were the main new line of attack, but they were not the only one.
This section records four other routes in enough detail to prevent the same hidden loss from
being rediscovered under different notation.

\subsection{Transferring convergents of \texorpdfstring{$\beta$}{beta} to
\texorpdfstring{$\vartheta$}{theta}}
\label{subsec:transfer}

Let $p/q$ be a continued-fraction convergent of a hypothetical primitive generator $\beta$ and
put
\[
 \varepsilon=q\beta-p.
\]
For all sufficiently large $q$, $|\varepsilon|<1/2$.  Define the two nonzero integer gaps
\begin{equation}\label{eq:D2D3}
 D_2=M^q-2^p=2^p(2^\varepsilon-1),
 \qquad
 D_3=A^q-3^p=3^p(3^\varepsilon-1).
\end{equation}
They manufacture an exact rational approximation to $\vartheta$:
\begin{equation}\label{eq:theta-transfer}
 R_{p,q}:=
 \frac{2^pD_3}{3^pD_2}
 =\frac{3^\varepsilon-1}{2^\varepsilon-1}
 \in\Q.
\end{equation}
Taylor expansion at zero gives
\begin{equation}\label{eq:theta-transfer-error}
 R_{p,q}=\vartheta+
 \frac{\log3\,(\log3-\log2)}{2\log2}\,\varepsilon
 +O(\varepsilon^2),
\end{equation}
so $|R_{p,q}-\vartheta|\ll|\varepsilon|$.

This looks promising because $\vartheta$ has an explicit irrationality measure.  Wu's
simultaneous linear-independence estimate for $1,\log2,\log3,\log5$ gives, after setting the
unused coefficients to zero, a conservative exponent
\begin{equation}\label{eq:wu-measure}
 \left|\vartheta-\frac rs\right|
 \ge c_\vartheta s^{-\mu_\vartheta},
 \qquad \mu_\vartheta=8.616
\end{equation}
for all sufficiently large denominators $s$ \cite{Wu2003}.  The obstacle is the denominator
created by \eqref{eq:theta-transfer}.  Before reduction it is at most
\begin{equation}\label{eq:transfer-denominator}
 3^p|D_2|
 \ll 3^p2^p|\varepsilon|
 =6^p|\varepsilon|.
\end{equation}
Combining \eqref{eq:theta-transfer-error}--\eqref{eq:transfer-denominator} with
\eqref{eq:wu-measure} gives only
\begin{equation}\label{eq:transfer-lower}
 |\varepsilon|
 \gg 6^{-p\mu_\vartheta/(1+\mu_\vartheta)}.
\end{equation}
Since
\[
 6^{-\mu_\vartheta/(1+\mu_\vartheta)}
 =2^{-2.316\ldots},
\]
this is substantially weaker than the elementary bound
$|\varepsilon|\gg2^{-p}$ obtained directly from $D_2\ne0$.  Sharper real irrationality
measures cannot repair the architecture: every irrationality exponent is at least two, and
the unreduced denominator already contains both $2^p$ and $3^p$.

\begin{proposition}[Exact denominator-loss diagnosis]\label{prop:transfer-loss}
The convergent-transfer construction can improve the elementary gap only if one proves an
exponential lower bound for
\begin{equation}\label{eq:transfer-gcd}
 G_{p,q}=\gcd(2^pD_3,3^pD_2)
\end{equation}
that survives reduction of $R_{p,q}$.  In the absence of such a bound, every known
irrationality measure for $\vartheta$ is defeated by the factor $6^p$ in
\eqref{eq:transfer-denominator}.
\end{proposition}

No mechanism forcing $G_{p,q}$ to be exponentially large was found.  The two gaps share the
same small real parameter $\varepsilon$, but real closeness does not impose a common prime
divisor.

\subsection{Two-row determinants and fixed-support \texorpdfstring{$S$}{S}-units}

For two block rows $i<j$, put $d=j-i$ and
$h=n_i-n_j=\lfloor j\beta\rfloor-\lfloor i\beta\rfloor$.  Their $2\times2$ determinant has
the exact factorization
\begin{equation}\label{eq:two-row-factor}
 m_{T,i}a_{T,j}-m_{T,j}a_{T,i}
 =6^{n_j}(MA)^i\bigl(2^hA^d-3^hM^d\bigr).
\end{equation}
Moreover
\begin{equation}\label{eq:two-row-ratio}
 \frac{2^hA^d}{3^hM^d}
 =\left(\frac32\right)^{\beta d-h}.
\end{equation}
Thus a good rational approximation $h/d$ to $\beta$ makes the bracket in
\eqref{eq:two-row-factor} a small \emph{relative} difference of two integers supported on the
fixed prime set dividing $6MA$.

Let
\[
 U=2^hA^d,
 \qquad V=3^hM^d,
 \qquad G=\gcd(U,V).
\]
The reduced nonzero integer $(U-V)/G$ gives
\begin{equation}\label{eq:sunit-gap}
 |\beta d-h|
 \gg_{M,A}\frac{G}{\max\{U,V\}}.
\end{equation}
The exponent of the right-hand side can be written explicitly as the positive part of a
valuation-vector difference:
\begin{equation}\label{eq:sunit-rate}
 \log\frac{U}{G}
 =\sum_p
 \max\!\left
 \{h\mathbf1_{p=2}+d v_p(A)
   -h\mathbf1_{p=3}-d v_p(M),0\right\}\log p.
\end{equation}
As $h/d\to\beta$, this is $\kappa_{M,A}d+o(d)$ for a positive constant
$\kappa_{M,A}$.  Hence \eqref{eq:sunit-gap} is merely exponential in $d$, much weaker than the
Baker--Matveev polynomial separation of Section~\ref{sec:baker}.

The $abc$ conjecture applied to the reduced equation $U/G-V/G=(U-V)/G$ would make the
difference nearly as large as a fixed positive power of $U/G$, but it would still allow a
relative gap $e^{-\epsilon d}$.  Continued-fraction gaps of order $1/d$ are much larger.
Thus this direct $S$-unit equation does not become a proof even after inserting $abc$; one
would need an additional theorem controlling the prime support or radical of the difference.

\subsection{Automaticity and the synchronized Sturmian word}

The coding $s_k=\lfloor(k+1)\alpha\rfloor-\lfloor k\alpha\rfloor$ from
Section~\ref{sec:sturmian} is aperiodic Sturmian.  A standard theorem in automatic-sequence
theory says that no aperiodic Sturmian word is automatic in any integer base
\cite{Cobham1972,AlloucheShallit2003}.  One quick obstruction is frequency: the frequency of
$1$ in $s$ is the irrational number $\alpha$, whereas an automatic sequence having a letter
frequency has rational frequency.

This observation blocks a tempting but invalid shortcut.  The rationality of
$U=2^\alpha$ does \emph{not} make the threshold itinerary of
$R\mapsto UR/2^{\mathbf1_{UR\ge2}}$ finite-state; its exact rational denominators grow with
time.  The same is true for $V$ in base three.  Therefore Cobham's theorem cannot be invoked
merely because the same word is produced by a dyadic and a triadic rational system.  A viable
automata proof would first have to manufacture genuine finite kernels in two multiplicatively
independent bases, and doing so is already a new arithmetic theorem about $U$ and $V$.

\subsection{Real and \texorpdfstring{$p$}{p}-adic logarithms}

At the real place the counterexample is exactly
\begin{equation}\label{eq:real-log-wall}
 \log M\,\log3-\log A\,\log2=0.
\end{equation}
At $p=2$ or $3$, after stripping valuation and torsion parts, the units among
$M,A,2,3$ admit $p$-adic logarithms.  The determinant entries then have unit parts whose
$p$-adic logarithms are bilinear forms in the row and column coordinates.  This is the natural
language for the cancellation cascade isolated in Section~\ref{sec:firstlayer}.

What is missing is a transfer principle from \eqref{eq:real-log-wall} to a congruence or rank
drop among those $p$-adic logarithms.  Real exponentiation by the transcendental number
$\beta$ does not define a compatible $p$-adic exponent, and the equality $M=2^\beta$ has no
$p$-adic meaning.  Known $p$-adic linear-form theorems control a nonzero form once its integer
coefficients and algebraic bases are specified; they do not create a $p$-adic form from a real
one.

\begin{target}[Real/$p$-adic compatibility]\label{target:real-padic}
Find an algebraic construction, valid for the integer pair $(M,A)$ alone, that converts
\eqref{eq:real-log-wall} into either
\begin{enumerate}
\item a nontrivial congruence among $p$-adic unit logarithms with precision $\gg T$, repeated
      over $\gg T^{3/2}$ fiber pairs; or
\item a rank drop in the tropical initial matrices persisting through $\gg T$ valuation
      layers.
\end{enumerate}
Without such a construction, ordinary $p$-adic logarithm bounds point in the opposite
direction: they show that individual congruences are only polylogarithmically deep.
\end{target}

\subsection{Transcendental-curve point counting}

The graph $y=x^\vartheta$ is definable in the real exponential field, so Pila--Wilkie type
results~\cite{PilaWilkie2006} give subpolynomial bounds for rational points outside algebraic pieces.  Those generic
bounds are far too weak here.  The curve already contains the infinite sequence
$(2^n,3^n)$, giving $\asymp\log H$ integer points of height at most $H$, while a hypothetical
second generator produces $\asymp(\log H)^2$ points globally.  A proof by point counting would
therefore need a sharp $O(\log H)$ theorem with a structural classification of the extremal
sequence, not merely $O(H^\epsilon)$.  Such a theorem would essentially be another form of the
rank-one conclusion sought by the conjecture.

\section{Exact sufficient conditions extracted from the new analysis}
\label{sec:sufficient}

A useful continuation should not stop at saying that the present determinant fails.  The
calculations above isolate numerical thresholds which, if crossed by a future theorem, would
settle the conjecture.  This section records those thresholds as precise implications.

\subsection{The mixed-block cancellation-excess criterion}

Let
\[
 \mathcal P=\{p:p\mid6MA\}
\]
be the fixed support of all determinant entries, and retain the mixed block determinant $D_T$,
its tropical minima $\tau_{p,T}$, and the archimedean scale $H_T$.  Define the support-prime
cancellation excess
\begin{equation}\label{eq:cancellation-excess}
 \mathcal E_T
 =\sum_{p\in\mathcal P}
   \bigl(v_p(D_T)-\tau_{p,T}\bigr)\log p
 \ge0
\end{equation}
and the exact tropical deficit
\begin{equation}\label{eq:tropical-deficit}
 \delta(M,A)=\frac38\sum_{p\in\mathcal P}\mathfrak G(c_p).
\end{equation}
Theorem~\ref{thm:mixed-gini} says
$\log Q_T^{\rm trop}=(1-\delta(M,A)+o(1))H_T$, while the elementary Leibniz upper bound is
$\log D_T\le H_T+o(H_T)$.

\begin{theorem}[Cancellation-excess criterion]\label{thm:excess-criterion}
If a hypothetical primitive counterexample satisfies
\begin{equation}\label{eq:excess-threshold}
 \liminf_{T\to\infty}\frac{\mathcal E_T}{H_T}
 >\delta(M,A),
\end{equation}
then no such counterexample exists.  In particular it is enough to prove an excess greater
than $H_T/5$ uniformly in all primitive branches; in the cross-coprime branch the exact
threshold is at least $3H_T/10$.
\end{theorem}

\begin{proof}
The selected support primes alone give
\[
 \log D_T
 \ge\sum_{p\in\mathcal P}v_p(D_T)\log p
 =\log Q_T^{\rm trop}+\mathcal E_T.
\]
Under \eqref{eq:excess-threshold}, the right-hand side exceeds $H_T$ by a fixed positive
multiple for all large $T$, contradicting the Leibniz upper bound.  The numerical thresholds
follow from Theorem~\ref{thm:four-fifths} and Corollary~\ref{cor:seven-tenths}.
\end{proof}

This criterion is not tautological: it involves only valuations at the fixed finite set of
primes already present in $6MA$, whereas $D_T$ can acquire arbitrary new prime factors.  It
states exactly how much same-support cancellation is required.

\subsection{A denominator-cancellation criterion}

For the transferred approximant $R_{p,q}$ of Section~\ref{subsec:transfer}, let $s_{p,q}$ be
its reduced denominator.  If one could prove along infinitely many convergents that
\begin{equation}\label{eq:poly-reduced-denominator}
 s_{p,q}\ll q^\eta
 \qquad\text{for some }\eta<\frac{1+\mu_\vartheta}{\mu_\vartheta},
\end{equation}
then Wu's measure would give
$|R_{p,q}-\vartheta|\gg q^{-\eta\mu_\vartheta}$, whereas
\eqref{eq:theta-transfer-error} and the convergent bound give
$|R_{p,q}-\vartheta|\ll q^{-1}$.  Since
$\eta\mu_\vartheta<1+\mu_\vartheta$ is not by itself enough, one tracks
$|\varepsilon|<1/q_{+}$ more precisely; a contradiction follows whenever
$q_{+}\gg q^{\eta\mu_\vartheta+\epsilon}$ along the same subsequence.  A simpler sufficient
condition is
\begin{equation}\label{eq:very-small-reduced-denominator}
 s_{p,q}=q^{o(1)}
\end{equation}
for infinitely many convergents with $q_{+}\ge q^{1+\epsilon}$.

Equivalently, the gcd $G_{p,q}$ in \eqref{eq:transfer-gcd} would have to cancel essentially the
entire exponential factor $6^p|\varepsilon|$, leaving a subpolynomial denominator.  This is a
sharp formulation of the otherwise vague request for ``correlation between the two integer
gaps.''

\subsection{An automaticity criterion}

The synchronized coding gives another exact, if presently inaccessible, implication.

\begin{proposition}[Finite-kernel criterion]\label{prop:automatic-criterion}
If under a hypothetical counterexample the common word
$s_k=\lfloor(k+1)\alpha\rfloor-\lfloor k\alpha\rfloor$ were automatic in any integer base
$q\ge2$, then the Alaoglu--Erd\H{o}s conjecture would follow.
\end{proposition}

\begin{proof}
An automatic Sturmian word is ultimately periodic.  Its slope, equal to the frequency of the
letter $1$, is therefore rational.  But $\alpha=\beta-\lfloor\beta\rfloor$ would then be
rational, and the rational-exponent theorem would make $\beta$ integral.
\end{proof}

The missing step is not symbolic dynamics; it is proving automaticity from the two rational
multipliers $U$ and $V$.  The growing denominators in \eqref{eq:normalized-orbits} show exactly why the
obvious state space is infinite.

\section{Branch-sensitive research program}
\label{sec:program}

The source report already classifies a hypothetical pair much more sharply than ``two
integers on a curve.''  The next attack should exploit that classification rather than seek a
single estimate that treats every branch identically.

\subsection{Branch I: independent external valuation directions}

Here the second-iterate kernel $K_\beta$ is zero, equivalently the external prime-valuation
vectors of $M$ and $A$ are independent.  Corollary~\ref{cor:gini-kernel} gives a strict
angular defect $\Delta_{\rm ext}>0$ in the all-prime tropical fraction.  That makes the
termwise divisor farther from the archimedean upper bound, but the same independence also
supplies at least two genuinely different external primes at which cancellation can be
studied.

The concrete target is a simultaneous $p$-adic rank theorem for the tropical initial matrices:
show that if cancellation persists through $L$ layers at two non-collinear external primes,
then an integer relation appears among the two external valuation vectors.  Such a theorem
would contradict $K_\beta=0$.  Existing one-prime linear-form bounds do not compare the
residue kernels at different primes; this is the genuinely new ingredient required.

\subsection{Branch II: a common external valuation ray}

When $K_\beta\ne0$, all external vectors are collinear and the source report gives the exact
M\"obius-orbit condition
\begin{equation}\label{eq:mobius-branch}
 \beta=\frac{u+v\vartheta}{r+s\vartheta}
 \qquad(u,v,r,s\in\Q,\\ us-vr\ne0).
\end{equation}
This is the branch in which the external triangle inequality can be an equality.  It is also
arithmetically lower-dimensional: after clearing denominators, the external parts of $M$ and
$A$ are powers of one common primitive core.

The best next step is therefore not a generic all-prime determinant.  It is a mixed
real/radical invariant built from the canonical core and the two binomial fields already
formalized in the repository.  Useful intermediate targets are:
\begin{enumerate}
\item compare the field norms of the canonical radicals on the two sides of
      \eqref{eq:mobius-branch};
\item show that simultaneous radical degrees force an additional common perfect-power degree,
      contradicting affine primitivity; or
\item prove that the M\"obius coefficients in \eqref{eq:mobius-branch} cannot have the sign
      pattern imposed by nonnegative prime valuations unless one output is $\{2,3\}$-smooth.
\end{enumerate}
The latest repository modules
\module{CanonicalRadicalFieldNorm}, \module{CanonicalRadicalValuation},
\module{CanonicalRadicalSharpness}, \module{CanonicalRadicalDivisibleHull}, and
\module{SimultaneousRadicalDegrees} provide much of the algebraic plumbing for this branch.

\subsection{Branch III: one smooth output}

If $M$ is $\{2,3\}$-smooth, then $\beta=a+d\vartheta$ with $d>0$; if $A$ is smooth, there is a
symmetric affine formula.  The source corpus proves that both outputs cannot be smooth and
that the other output contains a prime at least five.  In the extreme case $M=3^d$, the
problem is the localized radical assertion
\[
 (3^d)^\vartheta\notin\Z[1/3].
\]

This branch is the smallest transcendence-theoretic wall and deserves separate treatment.
Potentially useful objects are the irreducible binomial for the least localized radical, its
field norm, and reduction at primes over two and three.  A successful local argument must
remain sensitive to the real identity defining $\vartheta$; ordinary ideal factorization
alone is already exhausted by the canonical-radical modules.

\subsection{Branch IV: cross-coprime primitive outputs}

The cross-coprime branch has the strongest determinant constants and the cleanest residue
matrices.  Here the exact first tropical forms are the fiber Vandermondes of
Section~\ref{sec:firstlayer}.  This is the best laboratory for determining whether the
cancellation cascade is real or illusory.

The next computational experiment should not search for counterexamples.  For synthetic odd
pairs $(M,A)$ and Beatty depths generated by $\beta=\log_2M$, compute
\[
 v_2(D_T)-\tau_{2,T},
 \qquad
 v_3(D_T)-\tau_{3,T}
\]
for increasing $T$, together with the histogram of tropical layers contributing at the final
valuation.  Small exact experiments performed during this work show an excess compatible with
$N^{3/2}$ rather than $N^{5/2}$, matching Theorem~\ref{thm:first-layer-small}; they are
reconnaissance only, since a synthetic $A$ does not satisfy $A=3^\beta$.  A proof should seek a
uniform $o(H_T)$ bound, probably through $p$-adic exponential-polynomial zero estimates.

\subsection{Archimedean priority: solve the triangular HCIZ limit}

The complete grid finally supplies the missing $T^2$ rows, and its empirical measures are
explicit.  This makes Target~\ref{target:hciz} the highest-priority analytic problem.  The
required calculation is a large-$N$ asymptotic for the HCIZ integral with both spectra having
density $2t$ on $[0,1]$, under the coupled scaling $x_i y_j\asymp T^2$ and
$N\asymp T^2$.  Even a nonsharp upper bound improving the largest-term energy by more than
$1-\pi/4$ would materially change the $p$-adic comparison.

A practical sequence is:
\begin{enumerate}
\item numerically evaluate $\log\det(e^{x_i y_j})$ using high-precision total-positive
      algorithms for moderate $T$;
\item subtract the exact Vandermonde/factorial terms in \eqref{eq:hciz-det} to estimate the
      free-energy limit;
\item derive the corresponding variational problem using known HCIZ large-deviation methods;
\item only then reinsert the primewise valuation lower bounds.
\end{enumerate}
This route is attractive because the analytic limit does not depend on the unknown prime
factorization of $M$ and $A$; only $\beta$ enters the overall scaling.

\subsection{Finite verification priority}

The software floor is now $x>27$.  The immediate formal target is not another raw doubling of
the MPFR scan but a replayable certificate architecture.  A realistic hierarchy is:
\begin{enumerate}
\item a fast untrusted generator records only near-boundary pairs and rational interval data;
\item a small verified checker proves each interval enclosure and the monotone coverage between
      records;
\item chunks are combined by a theorem whose statement mentions only integer endpoints;
\item the final theorem is imported into the aggregate surface without
      \lean{native\_decide} or an opaque external axiom.
\end{enumerate}
The source-hash discrepancy found in Section~\ref{sec:finite27} makes this especially
worthwhile: mathematical replay is more robust than a manifest that identifies an external
binary only by prose.

\section{Lean formalization blueprint}
\label{sec:formalization}

The source corpus is already unusually strong on exact conditional structure.  The most useful
continuation is therefore not a second informal reimplementation of those results, but a
small collection of sharply separated modules whose statements expose the genuinely new
invariants.  This section gives a dependency order designed to keep arithmetic, asymptotics,
and transcendence inputs from leaking into one another.

\subsection{Proposed module graph}

\begin{longtable}{@{}>{\raggedright\arraybackslash}p{0.22\textwidth}>{\raggedright\arraybackslash}p{0.37\textwidth}>{\raggedright\arraybackslash}p{0.31\textwidth}@{}}
\toprule
Proposed module & Principal statement & Main dependencies \\
\midrule
\endfirsthead
\toprule
Proposed module & Principal statement & Main dependencies \\
\midrule
\endhead
\module{SturmianNormalization} & Exact recurrences for $R_k,S_k$, common jump word $s_k$, and the exponent/output correspondence. & Existing primitive-generator and denominator-normalization modules; elementary floor arithmetic. \\
\addlinespace
\module{ReducedOrbitDenominators} & Lowest-term denominators $q_k^{(2)},q_k^{(3)}$ and the full-depth alternative from $\min(v_2(M),v_3(A))=0$. & \module{SturmianNormalization}; unique factorization and valuations. \\
\addlinespace
\module{EffectiveGeneratorSeparation} & An abstract interface turning a two-logarithm lower bound into $|\beta-p/q|\ge c q^{-\mu}$ and a polynomial recurrence for continued-fraction denominators. & A stated external Baker--Matveev hypothesis, plus existing continued-fraction lemmas. \\
\addlinespace
\module{BlockVandermondeValuation} & Exact $v_2$ and fixed-support $v_p$ formulas for the conditional block Vandermonde; output analogue. & Existing exact block parametrization; valuation of a difference with unequal valuations. \\
\addlinespace
\module{OrdinaryAdelicFraction} & The asymptotic fraction \eqref{eq:ordinary-fraction} and the universal $1-\log2/(3\log6)$ ceiling. & \module{BlockVandermondeValuation}; polynomial floor sums; primitive-output depth theorem. \\
\addlinespace
\module{TriangleGini} & $\mathfrak G$ is a norm and satisfies the explicit formula \eqref{eq:G-explicit}. & Elementary measure integration, or a finite polygonal integration library. \\
\addlinespace
\module{FiniteTropicalAssignment} & Exact rearrangement formula for a finite row list and a finite column-weight list. & Sorting/permutations; finite rearrangement inequality. \\
\addlinespace
\module{MixedColumnRiemannSums} & Lattice count \eqref{eq:semigroup-count} and convergence of scaled low-weight columns to uniform measure on $\DeltaTri$. & Two-dimensional lattice-point estimates and Riemann sums. \\
\addlinespace
\module{MixedAdelicGini} & The limit \eqref{eq:mixed-gini-limit}, conservation $\sum_pc_p=0$, and branch-sensitive three-vector bound. & Previous three modules; finite prime support. \\
\addlinespace
\module{MixedTropicalCeilings} & Universal $4/5$, cross-coprime $7/10$, and structural $3/10$ constants. & \module{MixedAdelicGini}; primitive-output depth theorem; explicit values from \module{TriangleGini}. \\
\addlinespace
\module{FullGridTransport} & Row/column empirical distributions, countermonotone energy $\pi/8$, maximum energy $1/2$, and the ratio $\pi/4$. & Triangular lattice sums, quantile integration, rearrangement. \\
\addlinespace
\module{TropicalInitialFibers} & Factorization of the first $2$- and $3$-adic initial forms into ordinary Vandermondes on column fibers. & Determinants over valued rings; stable sorting by valuations. \\
\addlinespace
\module{InitialFiberValuationBound} & Under a standard $p$-adic logarithm hypothesis, the first tied layer contributes $o(H_T)$. & \module{TropicalInitialFibers}; an abstract two-logarithm valuation bound. \\
\addlinespace
\module{FiniteCheck2Pow27} & A kernel-replayable certificate excluding non-powers of two below $2^{27}$. & Existing finite-check certificate architecture; generated chunk certificates. \\
\bottomrule
\end{longtable}

The first two modules are almost entirely combinatorial.  The determinant modules should not
import transcendence theory.  Conversely, the Baker and $p$-adic logarithm interfaces should
be abstract propositions whose classical realization is documented outside the kernel unless
and until Mathlib contains a suitable theorem.  This makes the trust boundary explicit.

\subsection{Finite statements before asymptotics}

The continuum notation in Sections~\ref{sec:mixed} and~\ref{sec:fullgrid} is compact, but Lean
should first see finite inequalities.  For rows $0\le k<N$ and columns $z_j\in\R$, the finite
rearrangement lemma should assert
\begin{equation}\label{eq:finite-rearrangement}
  \min_{\sigma\in S_N}\sum_{k=0}^{N-1} k z_{\sigma(k)}
  =\sum_{k=0}^{N-1} k z_{N-1-k}^{\uparrow},
\end{equation}
where $z^{\uparrow}$ is the nondecreasing rearrangement.  Every tropical assignment in this
report is an affine rescaling of \eqref{eq:finite-rearrangement}.  The limiting Gini identity
can then be obtained by three auditable steps:
\begin{enumerate}
\item prove convergence of empirical distributions for the triangular lattice;
\item prove convergence of the first moments and pairwise absolute-difference moments;
\item invoke the finite identity
\[
 \frac1{N^2}\sum_{i,j}|z_i-z_j|
 =\frac2{N^2}\sum_{j=0}^{N-1}(2j-N+1)z_j^{\uparrow}.
\]
\end{enumerate}
This route avoids formalizing optimal transport as a general theory.

\subsection{Exact theorem sketches}

The intended public statements can be phrased close to the mathematics.

\begin{lstlisting}[style=leanstyle]
/-- A primitive counterexample determines one common lower mechanical word
    driving both normalized rational orbits. -/
theorem exists_synchronized_sturmian_data
    (hfail : ExistsNonintegralSolution) :
    Exists fun B : Nat => Exists fun alpha U V : Real =>
      0 < alpha /\ alpha < 1 /\ Irrational alpha /\
      U = 2 ^ alpha /\ V = 3 ^ alpha /\
      (forall k : Nat,
        let e := Int.floor (k * alpha)
        normalizedTwo k = U ^ k / 2 ^ e /\
        normalizedThree k = V ^ k / 3 ^ e) := by
  ...

/-- The automatic all-prime divisor of the balanced mixed determinant
    cannot fill more than four fifths of its leading height. -/
theorem mixed_tropical_divisor_limsup_le_four_fifths
    (hprimitive : min (vTwo M) (vThree A) = 0) :
    limsup (fun T => log (tropicalDivisor T) / leadingHeight T)
      <= (4 : Real) / 5 := by
  ...

/-- The complete conditional grid has countermonotone/maximal transport
    ratio pi/4. -/
theorem full_grid_transport_ratio :
    countermonotoneEnergy / maximalEnergy = Real.pi / 4 := by
  ...
\end{lstlisting}

The actual APIs will need natural-valued denominator exponents rather than casts of negative
integers, and the asymptotic statements will likely use \lean{Tendsto} and \lean{IsLittleO}
rather than a custom \lean{limsup}.  The sketches are included to fix the mathematical
interfaces, not Lean syntax.

\subsection{Formalizing the first tied layer}

A useful reusable notion is the initial form of a determinant at a prime.  For a matrix over
$\Z$ and an assignment minimum $\tau_p$, reduce
\[
 p^{-\tau_p}\det(E_{ij})\pmod p.
\]
Only minimizing permutations survive.  In the cross-coprime branch the minimizing
permutations preserve equal-$u$ fibers at $p=2$ and equal-$v$ fibers at $p=3$; the surviving
sum is a product of smaller determinants.  This should be proved as a purely finite
combinatorial determinant identity.  The later valuation estimate can then be stated
separately:
\begin{equation}\label{eq:initial-form-bound-formal}
  v_2(I_{2,T})+v_3(I_{3,T})
  =O\!\left(N^{3/2}(\log N)^2\right).
\end{equation}
The big-$O$ exponent is not sacred; any $o(T^{5/2})$ theorem is sufficient for the qualitative
conclusion.

\subsection{Certificate architecture for the new finite bound}

The existing verified bound through $4096$ should remain the trusted baseline until a new
certificate is accepted.  A practical route to $2^{27}$ is hierarchical:
\begin{enumerate}
\item partition $[1,2^{27})$ into fixed-width chunks;
\item for each non-power-of-two $m$, store an integer nearest-value locator $n_m$ and rational
      lower/upper enclosures proving either $m^\vartheta<n_m$ or $m^\vartheta>n_m$;
\item hash each chunk and prove a small checker correct in Lean;
\item combine the chunk theorems into one interval theorem.
\end{enumerate}
The MPFR run in Section~\ref{sec:finite27} is reconnaissance and a source of extremal cases;
it is not the certificate itself.

\section{Numerical sanity checks and reproducibility}
\label{sec:numerics}

None of the limiting constants in the determinant theorems depends on numerical evidence.
The computations below were used to catch normalization errors, check convergence rates, and
stress-test the closed forms before writing the proofs.

\subsection{Balanced-block assignment constants}

For each $N$, take the first $N$ pairs $(u,v)$ ordered by $u+\vartheta v$, pair them against
$N$ equally spaced row parameters, and solve the one-dimensional assignment problems by
sorting.  In the cross-coprime model, the structural places should tend to $3/10$ and the
minimal all-prime three-vector model to $7/10$.

\begin{table}[htbp]
\centering
\small
\begin{tabular}{@{}rcc@{}}
\toprule
$N$ & structural $2,3$ fraction & all-prime three-vector fraction \\
\midrule
25   & 0.3224571184 & 0.7258906037 \\
50   & 0.3135103379 & 0.7141968856 \\
100  & 0.3050363664 & 0.7053583986 \\
250  & 0.3024422068 & 0.7025311250 \\
500  & 0.3012614628 & 0.7013136826 \\
1000 & 0.3006780555 & 0.7006854844 \\
2500 & 0.3002789586 & 0.7002894348 \\
5000 & 0.3001354539 & 0.7001384826 \\
\midrule
Limit & $3/10$ & $7/10$ \\
\bottomrule
\end{tabular}
\caption{Finite assignment sums converging to the exact mixed-block constants.}
\label{tab:mixed-numerics}
\end{table}

The convergence from above is consistent with the $O(L^{-1})=O(N^{-1/2})$ boundary error in
triangular lattice counting.  The table also exposes the error in the proposed $2/3$ ceiling:
the finite all-prime ratio is already near $0.70$, and the exact limit is $7/10$.

\subsection{Complete-grid transport ratio}

For a representative hypothetical height parameter $\beta=27.314159$, form the row lattice
$n+\beta k\le T$ and the matching low-weight column lattice.  Normalize the
countermonotone assignment by the largest-term assignment.  The finite ratios are:

\begin{table}[htbp]
\centering
\small
\begin{tabular}{@{}rrc@{}}
\toprule
$T$ & number of rows & transport ratio \\
\midrule
50  & 73    & 0.7415113 \\
100 & 237   & 0.7559915 \\
200 & 841   & 0.7671191 \\
400 & 3157  & 0.7750371 \\
800 & 12210 & 0.7796781 \\
\midrule
Limit & --- & $\pi/4=0.7853981634\ldots$ \\
\bottomrule
\end{tabular}
\caption{Finite full-grid countermonotone/maximal transport ratios.}
\label{tab:grid-numerics}
\end{table}

The slow boundary convergence is expected: both empirical triangles have $O(T)$ boundary
points against $O(T^2)$ bulk points, and the matching weight scale also grows with $T$.

\subsection{Directed-rounding scan manifest}

The scan in Section~\ref{sec:finite27} used the literal C listing extracted from the attached
LaTeX source.  Its SHA-256 digest is
\begin{center}
\sha{f888c0d89be9366f01e2d26f339dcbe446556d39d4c5dbcf8f579ec0d0052e1d}.
\end{center}
This differs from the digest printed in the old report.  To ensure the discrepancy was only a
manifest error, the extracted source was first rerun on the old overlap interval
$[29{,}360{,}128,33{,}554{,}432)$; it reproduced the old extremal records and reported zero
failures.

The extension interval $[2^{26},2^{27})$ was split into sixteen chunks of width $2^{22}$ and
run at $192$ bits.  The concatenated chunk log has SHA-256
\begin{center}
\sha{b11843448cd8946442aba13758b5b666e4ac6eae50d2243fbdc32d83274cea8d}.
\end{center}
The two new extremal records were replayed at $320$ and $512$ bits.  A representative build
and invocation sequence was:
\begin{lstlisting}[style=pythonstyle,language=bash]
cc -O3 -std=c17 mpfr_scan.c \
  /lib/x86_64-linux-gnu/libmpfr.so.6 -lgmp -lm -o mpfr_scan

./mpfr_scan 67108864 71303168 192
# repeat consecutive width-2^22 chunks through 134217728
sha256sum extension-through-2pow27.log
\end{lstlisting}
The scanner excludes powers of two from the tested population because they are the known
solutions.  Across the extension it checked $67{,}108{,}863$ non-powers of two and reported
zero failures.

\subsection{Extremal replay records}

The narrowest new lower-directed separation was
\[
 m=132{,}833{,}765,
 \qquad
 n=7{,}501{,}348{,}219{,}569,
 \qquad
 \text{directed margin }5.9082259992879571\times10^{-21}.
\]
The corresponding high-precision ordinary difference is approximately
\[
 m^\vartheta-n=6.39397545334704\times10^{-8}.
\]
The narrowest new upper-directed separation was
\[
 m=129{,}799{,}345,
 \qquad
 n=7{,}231{,}571{,}292{,}851,
 \qquad
 \text{directed margin }8.3073231820333117\times10^{-22},
\]
with
\[
 n+1-m^\vartheta=8.666990435581896\times10^{-9}.
\]
These decimal differences are explanatory only.  The certification decision uses the MPFR
outward-rounded interval endpoints and the scanner's locator checks.

\section{Status matrix and dependency audit}
\label{sec:status}

\begin{longtable}{@{}>{\raggedright\arraybackslash}p{0.37\textwidth}>{\raggedright\arraybackslash}p{0.18\textwidth}>{\raggedright\arraybackslash}p{0.35\textwidth}@{}}
\toprule
Statement & Status & Dependency or caveat \\
\midrule
\endfirsthead
\toprule
Statement & Status & Dependency or caveat \\
\midrule
\endhead
Exact primitive monoid $\Sol=\Nzero+\Nzero\beta$ and candidate monoid & \statuslean & Inherited from the attached report. \\
\addlinespace
Primitive-depth condition $\min(v_2(M),v_3(A))=0$ & \statuslean & Inherited; central to both universal ceilings. \\
\addlinespace
Synchronized dyadic--triadic Sturmian recurrence & \statusnew & Elementary from the exact monoid and floor normalization. \\
\addlinespace
Lowest-term denominator formulas and full-depth alternative & \statusnew & Unique factorization plus primitive depth. \\
\addlinespace
Effective polynomial irrationality measure for fixed $M$ & \statusclassical{} plus \statusnew & Baker--W"ustholz/Matveev; constants depend effectively on $M$. \\
\addlinespace
Polynomial recurrence of continued-fraction denominators & \statusnew & Standard convergent inequality applied to the previous row. \\
\addlinespace
Exact ordinary block valuations at all input-support primes & \statusnew & Unequal valuation of the two monomials in every difference. \\
\addlinespace
Ordinary product fraction \eqref{eq:ordinary-fraction} & \statusnew & Exact floor sums; output endpoint errors are $O(T^2)$. \\
\addlinespace
Universal ordinary ceiling $0.871049\ldots$ & \statusnew & Primitive depth only; not a bound on cancellation from new primes in differences. \\
\addlinespace
Explicit triangular Gini norm & \statusnew & Direct integration. \\
\addlinespace
Mixed all-prime tropical limit \eqref{eq:mixed-gini-limit} & \statusnew & Rearrangement plus triangular Riemann sums. \\
\addlinespace
Universal $4/5$ mixed ceiling & \statusnew & Conservation, Gini triangle inequality, primitive depth. \\
\addlinespace
Cross-coprime $7/10$ and structural $3/10$ & \statusnew & Specialization $a=b=c=d=0$. \\
\addlinespace
Full-grid $\pi/4$ ceiling & \statusnew & Countermonotone and comonotone transport for the common quantile $\sqrt t$. \\
\addlinespace
HCIZ determinant identity & \statusclassical & Exact finite identity; the needed large-$N$ inequality remains open. \\
\addlinespace
First tied-layer fiber factorization & \statusnew & Exact initial-form determinant algebra in the cross-coprime branch. \\
\addlinespace
First tied layer is $o(H_T)$ & \statusclassical{} plus \statusnew & Standard $p$-adic two-logarithm estimates; no claim about higher valuation layers. \\
\addlinespace
Mixed cancellation-excess criterion & \statusnew & Immediate sufficient condition from integrality and the Gini deficit. \\
\addlinespace
Software exclusion through $2^{27}$ & \statuscomp & Complete MPFR directed-rounding scan; not a Lean proof. \\
\addlinespace
Kernel exclusion through $4096$ & \statuslean & Remains the strongest formal finite theorem audited in the source report. \\
\addlinespace
The Alaoglu--Erd\H{o}s conjecture itself & \statusopen & Equivalent to the unresolved special four-exponentials determinant. \\
\bottomrule
\end{longtable}

\subsection{Logical dependency graph}

The new deductions can be summarized as
\[
\begin{array}{c}
\text{exact primitive monoid}\ +\ \text{primitive depth}
\\[2pt]\Downarrow\\[-2pt]
\text{synchronized Sturmian system}
\qquad\text{and}\qquad
\text{exact block valuation profiles}
\\[2pt]\Downarrow\\[-2pt]
\begin{array}{c}
\text{ordinary all-prime fraction}\\[2pt]
\text{mixed Gini conservation law}\\[2pt]
\text{full-grid transport problem}
\end{array}
\\[2pt]\Downarrow\\[-2pt]
\begin{array}{c}
0.871049\ldots\text{ ordinary ceiling},\quad
4/5\text{ mixed ceiling},\quad
\pi/4\text{ full-grid ceiling}
\end{array}
\\[2pt]\Downarrow\\[-2pt]
\text{only cancellation beyond tropical minima or a sharper HCIZ bound can finish these routes.}
\end{array}
\]
Baker--Matveev separation and the finite scan are parallel inputs: they sharpen the possible
counterexample but do not enter the determinant ceiling proofs.

\section{Conclusions}
\label{sec:conclusion}

The Alaoglu--Erd\H{o}s conjecture remains unproved.  The continuation nevertheless changes the
shape of the search in several concrete ways.

First, the exact conditional monoid has a compact symbolic normal form.  After removing the
integer part of the primitive generator, one irrational Sturmian word controls both the
binary and ternary denominator jumps of two rational recurrences.  This exposes a possible
route through finite-state rigidity, but it also prevents a false shortcut: rational
multipliers do not by themselves make the word automatic.

Second, ordinary transcendence theory gives more Diophantine separation than the source
report's elementary argument records.  A fixed hypothetical generator is a ratio of two
logarithms of integers, so Baker--W"ustholz/Matveev yields an effective polynomial
irrationality measure.  Continued-fraction denominators therefore grow at most polynomially
from one stage to the next.  This is useful quantitative structure, but it does not conflict
with the $O(T)$ exact population of a dyadic block.

Third, the most direct $p$-adic determinant plan can now be audited with exact constants.
For ordinary block Vandermondes, every valuation at every prime already supporting the inputs
is accounted for, and primitive depth leaves at least $12.8950\ldots\%$ of the leading
archimedean budget uncovered.  For the balanced mixed determinant, all primes collapse into a
finite family of slope vectors summing to zero.  The triangular Gini norm then proves that
entrywise divisibility misses at least $20\%$ universally, at least $30\%$ in the
cross-coprime branch, and at least $70\%$ if one uses only the structural places $2,3$ there.
The complete two-dimensional grid improves the geometry but still leaves a termwise ceiling
of $\pi/4$.

These are no-go theorems for particular mechanisms, not evidence that the conjecture is false.
They identify the exact remaining mechanisms:
\begin{enumerate}
\item a leading-order cascade of cancellations among many equal-valuation Leibniz terms;
\item an archimedean determinant estimate substantially smaller than its largest-term bound,
      most naturally through the HCIZ free energy;
\item a genuinely mixed real/$p$-adic construction that turns
      $\log M\log3=\log A\log2$ into deep congruences; or
\item a rigidity theorem for the synchronized rational Sturmian system.
\end{enumerate}
The first cancellation layer is provably too small, so any successful $p$-adic proof must
propagate through many valuation layers rather than extracting one convenient Vandermonde
factor.

Finally, the directed-rounding computation was extended by one full binary exponent.  At the
software-trust level, no non-power-of-two candidate occurs below $2^{27}$, so a counterexample
would satisfy $x>27$.  The formal Lean boundary remains $x\ge12$ until a generated certificate
is checked by the kernel.  The old scanner-source hash printed in the attached report is
incorrect; the literal source, overlap replay, new log hash, and high-precision extremal
replays are recorded explicitly in this document.

\begin{resultbox}
\textbf{Bottom line.}
The conjecture is still open, but the continuation produces new exact mathematics rather than
only a list of ideas: a synchronized Sturmian equivalence; effective polynomial separation;
closed all-prime Vandermonde asymptotics; a triangular Gini norm governing mixed tropical
divisibility; universal $4/5$, $7/10$, $3/10$, and $\pi/4$ ceilings; a first-layer
cancellation theorem; exact sufficient thresholds for future work; and a reproducible
software exclusion through $2^{27}$.  The next decisive calculation is the triangular HCIZ
free energy or, on the arithmetic side, a proof of a genuinely leading tropical cancellation
cascade.
\end{resultbox}

\begin{thebibliography}{99}

\bibitem{AlaogluErdos1944}
L.~Alaoglu and P.~Erd\H{o}s,
\newblock \emph{On highly composite and similar numbers},
\newblock Transactions of the American Mathematical Society \textbf{56} (1944), 448--469.
\newblock DOI: \href{https://doi.org/10.1090/S0002-9947-1944-0011087-2}{10.1090/S0002-9947-1944-0011087-2}.

\bibitem{Baker1975}
A.~Baker,
\newblock \emph{Transcendental Number Theory},
\newblock Cambridge University Press, 1975.

\bibitem{BakerWustholz1993}
A.~Baker and G.~W\"ustholz,
\newblock \emph{Logarithmic forms and group varieties},
\newblock Journal f\"ur die reine und angewandte Mathematik \textbf{442} (1993), 19--62.
\newblock DOI: \href{https://doi.org/10.1515/crll.1993.442.19}{10.1515/crll.1993.442.19}.

\bibitem{Matveev2000}
E.~M.~Matveev,
\newblock \emph{An explicit lower bound for a homogeneous rational linear form in the logarithms of algebraic numbers. II},
\newblock Izvestiya: Mathematics \textbf{64}(6) (2000), 1217--1269.
\newblock DOI: \href{https://doi.org/10.1070/IM2000v064n06ABEH000314}{10.1070/IM2000v064n06ABEH000314}.

\bibitem{BugeaudLaurent1996}
Y.~Bugeaud and M.~Laurent,
\newblock \emph{Minoration effective de la distance $p$-adique entre puissances de nombres alg\'ebriques},
\newblock Journal of Number Theory \textbf{61}(2) (1996), 311--342.
\newblock DOI: \href{https://doi.org/10.1006/jnth.1996.0152}{10.1006/jnth.1996.0152}.

\bibitem{Bugeaud2002}
Y.~Bugeaud,
\newblock \emph{Linear forms in two $m$-adic logarithms and applications to Diophantine problems},
\newblock Compositio Mathematica \textbf{132}(2) (2002), 137--158.
\newblock DOI: \href{https://doi.org/10.1023/A:1015825809661}{10.1023/A:1015825809661}.

\bibitem{Yu1998}
K.~Yu,
\newblock \emph{$p$-adic logarithmic forms and group varieties I},
\newblock Journal f\"ur die reine und angewandte Mathematik \textbf{502} (1998), 29--92.
\newblock DOI: \href{https://doi.org/10.1515/crll.1998.090}{10.1515/crll.1998.090}.

\bibitem{Wu2003}
Q.~Wu,
\newblock \emph{On the linear independence measure of logarithms of rational numbers},
\newblock Mathematics of Computation \textbf{72}(242) (2003), 901--911.
\newblock DOI: \href{https://doi.org/10.1090/S0025-5718-02-01442-4}{10.1090/S0025-5718-02-01442-4}.

\bibitem{MorseHedlund1940}
M.~Morse and G.~A.~Hedlund,
\newblock \emph{Symbolic dynamics II. Sturmian trajectories},
\newblock American Journal of Mathematics \textbf{62}(1) (1940), 1--42.
\newblock DOI: \href{https://doi.org/10.2307/2371431}{10.2307/2371431}.

\bibitem{Lothaire2002}
M.~Lothaire,
\newblock \emph{Algebraic Combinatorics on Words},
\newblock Encyclopedia of Mathematics and its Applications 90, Cambridge University Press, 2002.

\bibitem{Cobham1972}
A.~Cobham,
\newblock \emph{Uniform tag sequences},
\newblock Mathematical Systems Theory \textbf{6} (1972), 164--192.
\newblock DOI: \href{https://doi.org/10.1007/BF01706087}{10.1007/BF01706087}.

\bibitem{AlloucheShallit2003}
J.-P.~Allouche and J.~Shallit,
\newblock \emph{Automatic Sequences: Theory, Applications, Generalizations},
\newblock Cambridge University Press, 2003.

\bibitem{HarishChandra1957}
Harish-Chandra,
\newblock \emph{Differential operators on a semisimple Lie algebra},
\newblock American Journal of Mathematics \textbf{79}(1) (1957), 87--120.
\newblock DOI: \href{https://doi.org/10.2307/2372387}{10.2307/2372387}.

\bibitem{ItzyksonZuber1980}
C.~Itzykson and J.-B.~Zuber,
\newblock \emph{The planar approximation. II},
\newblock Journal of Mathematical Physics \textbf{21}(3) (1980), 411--421.
\newblock DOI: \href{https://doi.org/10.1063/1.524438}{10.1063/1.524438}.

\bibitem{McSwiggen2018}
C.~McSwiggen,
\newblock \emph{The Harish--Chandra integral: an introduction with examples},
\newblock L'Enseignement Math\'ematique \textbf{67}(3--4) (2021), 229--299;
preprint arXiv:1806.11155.
\newblock DOI: \href{https://doi.org/10.4171/LEM/1017}{10.4171/LEM/1017}.

\bibitem{PilaWilkie2006}
J.~Pila and A.~J.~Wilkie,
\newblock \emph{The rational points of a definable set},
\newblock Duke Mathematical Journal \textbf{133}(3) (2006), 591--616.
\newblock DOI: \href{https://doi.org/10.1215/S0012-7094-06-13336-7}{10.1215/S0012-7094-06-13336-7}.

\bibitem{Waldschmidt2000}
M.~Waldschmidt,
\newblock \emph{Diophantine Approximation on Linear Algebraic Groups: Transcendence Properties of the Exponential Function in Several Variables},
\newblock Grundlehren der mathematischen Wissenschaften 326, Springer, 2000.

\bibitem{Roy1992}
D.~Roy,
\newblock \emph{Matrices whose coefficients are linear forms in logarithms},
\newblock Journal of Number Theory \textbf{41}(1) (1992), 22--47.
\newblock DOI: \href{https://doi.org/10.1016/0022-314X(92)90081-Y}{10.1016/0022-314X(92)90081-Y}.

\bibitem{ProveIt}
V.~Reshetnikov,
\newblock \emph{ProveIt repository},
\newblock \url{https://github.com/VladimirReshetnikov/ProveIt},
\newblock latest snapshot consulted for this report:
\sha{ff9a9f5aa3a3f02599de93fc053e60918f269855}.

\bibitem{UnifiedReport}
The ProveIt project,
\newblock \emph{The Alaoglu--Erd\H{o}s Integer-Exponent Conjecture: Unified Report},
\newblock attached LaTeX source \path{alaoglu-erdos-unified-report.tex}, August 2026.

\end{thebibliography}

\end{document}
